\documentclass[a4paper,fleqn]{ltxdoc}
\usepackage[margin=1in]{geometry}
%\input{pgfmanual-dvipdfm.cfg}
%\input{../../text-en/pgfmanual-en-main-preamble}
\usepackage[version=latest]{pgf}
\usepackage{xkeyval,calc,listings,tikz,fp,amsmath,amssymb}
\usepackage[T1]{fontenc}% big thanks to samcarter!
\usepackage{makeidx}
\makeindex
\usepackage{hyperref}
\hypersetup{%
        colorlinks=true,
        linkcolor=blue,
        filecolor=blue,
        urlcolor=blue,
        citecolor=blue,
        pdfborder=0 0 0,
}
\makeatletter          % see https://tex.stackexchange.com/q/33946
\input{pgfmanual.code} % 
\makeatother           % 
\input{pgfmanual-en-macros.tex} % link from
\makeatletter% recent LaTeX kernels list TAB in \dospecials, so codeexample
% would keep leading TABs as characters; restore the old behaviour
\def\code@catcode@hook{\catcode`\^^I=10\relax}%
\makeatother
% /usr/local/texlive/2020/texmf-dist/doc/generic/pgf/macros/pgfmanual-en-macros.tex
% or the equivalent on your installation
\def\pgfautoxrefs{1}
\usetikzlibrary{3dtools,3dtools.circleofsphere,calc,angles}
\DeclareMathOperator{\sd}{sd}
\begin{document}
\section{3D Tools}
\begin{tikzlibrary}{3dtools}
	This library provides additional tools to create 3d-like pictures. It is a
	collection of reasonably working tools, which however is not streamlined,
	and may be subject to substantial changes if the library ever happens to get
	further developed or published. This library loads the |3d|, |decorations|
	and |fpu| libraries.
\end{tikzlibrary}

\tikzname\ has the |3d| and |perspective| libraries which deal with the
projections of three-dimensional drawings. In addition there exist excellent
packages like \href{https://ctan.org/pkg/tikz-3dplot?lang=en}{|tikz-3dplot|}.
The purpose of this library is to provide some means to manipulate the
coordinates. It supports linear combinations of vectors, vector products and
scalar products. It also supports ``physical'' coordinates and means to
transform coordinates from one orthonormal basis to another orthonormal basis. 

\medskip

\noindent\textbf{Note:} Hopefully this library is only temporary and its
contents will be absorbed in slightly extended versions of the |3d| and |calc|
libraries.\footnote{Note that an earlier version of the library had been
uploaded to CTAN by the maintainer of \tikzname\ at that time. The maintainer
did not notify the author of the library, let alone ask for permission. Since
the library was in an even worse shape at that time, it got removed by the
author. For further information on this incident please consider contacting 
CTAN (rather than simply believing in chat messages, which are often not too
accurate).} The cleanest way will be to record a screen depth when ``saving'' a
coordinate with \tikzname. Some limited support of such a functionality is
provided in this library, see section~\ref{sec:PhysicalCoordinates}. However, a
full-fledged realization would require changes at the level of |tikz.code.tex|
and can only be done consistently if the maintainer(s) of \tikzname\ make some
fundamental changes. Even if this happens, dealing with node anchors and so on,
which so far are intrinsically two-dimensional (see, however,
\href{https://github.com/ZhiyuanLck/dnnplot}{https://github.com/ZhiyuanLck/dnnplot}
for an impressive counterexample) will be nontrivial because one will have to
specify the orientation in 3d as far as supported. Note also that it is quite
conceivable that the viewers in the future will be able to achieve 3d ordering,
so, in a way, recording the screen depth (see the |screendepth| function below)
will become almost mandatory at a given point. 

\medskip

\noindent\textbf{Filing a bug report or placing a feature request.} This library
is currently hosted under
\href{https://github.com/marmotghost/tikz-3dtools}{\texttt{https://github.com/marmotghost/tikz-3dtools}}.
% The author is also active at the \emph{noncommercial} Q \& A site
% \href{https://topanswers.xyz/tex}{\texttt{https://topanswers.xyz/tex}}, which
% offers, apart from the possibility of asking questions and browsing through the
% posts a chat in which one can discuss problems. 

\medskip

\noindent\textbf{Why is this library not on CTAN?} First of all, ideally this
library is only temporary. Secondly, it is also under construction. There are
many things that users may want to be able to use, such as spherical
coordinates, and which are not supported at all. Adding support for such
features my require syntax changes, thus making currently working code invalid.
If this is the case, it would not be helpful if the library had already been
made ``official''.

\medskip

\noindent\textbf{Code examples.} This manual contains a number of code examples.
They do not consist of complete \LaTeX\ documents but only excerpts. Almost all of
them require the |3dtools| library but they may require additional libraries that
do not get loaded by the |3dtools| library. All libraries that are required to
compile a code example are indicated. Each code example can be embedded in a minimal
example, which can have the form 
\[\begin{minipage}{0.8\textwidth}
\begin{tabular}{lp{1em}l}
|\documentclass[tikz]{standalone}| & & also loads \tikzname \\
|\usetikzlibrary{3dtools,|\dots|}| & & libraries as indicated\\
|\begin{document}| & & \\
\meta{code} & & code from the example\\
|\end{document}| & & \\
\end{tabular}
\end{minipage}
\]
It should be also possible to embed the code in generic \LaTeX\ documents as long as
in the preamble \tikzname\ gets loaded (e.g.\ with |\usepackage{tikz}|)  and the
relevant libraries get also loaded. Note that although above |ConTeXt| gets
mentioned this library has not been tested with |ConTeXt|. The only reason why
|ConTeXt| gets mentioned is that this gets automatically added by
|pgfmanual.code.tex|, which gets loaded to prepare this document.

\subsection{Coordinate computations}
\label{sec:3DCoordinateComputations}

The |3dtools| library has some options and styles for coordinate computations
and manipulations. 
\begin{key}{/tikz/3d parse}
        Parses an expression and inserts the result in form of a coordinate.
\end{key}
\begin{key}{/tikz/3d coordinate}
        Allow one to define a 3d coordinate from other coordinates.
\end{key}
Both keys support both symbolic and explicit coordinates.

\begin{codeexample}[width=6cm,preamble={\usetikzlibrary{3dtools}}]
\begin{tikzpicture}
 \path (1,2,3) coordinate (A) 
  (2,3,-1) coordinate (B) 
  (-1,-2,1) coordinate (C)
  [3d parse={0.25*(1,2,3)x(B)}] 
  	coordinate(D)
  [3d parse={0.25*(C)x(B)}] 
  	coordinate(E);
 \path foreach \X in {A,...,E} 
 {(\X) node[fill,inner sep=1pt,
 label=above:$\X$]{}};
\end{tikzpicture}
\end{codeexample}

Notice that, as of now, only the syntax |\path (1,2,3) coordinate (A);| works,
yet the syntax |\coordinate (A) at (1,2,3);| does \emph{not} work, but leads to
error messages.

\begin{codeexample}[width=6cm,preamble={\usetikzlibrary{3dtools}}]
\begin{tikzpicture}
 \path (1,2,3) coordinate (A) 
  (2,3,-1) coordinate (B) 
  (-1,-2,1) coordinate (C)
  [3d coordinate={(D)=0.25*(1,2,3)x(B)},
  3d coordinate={(E)=0.25*(C)x(B)},
  3d coordinate={(F)=(A)-(B)},];
 \path foreach \X in {A,...,E} 
 {(\X) node[fill,inner sep=1pt,
 label=above:$\X$]{}};
\end{tikzpicture}
\end{codeexample}

The actual parsing are done by the function |\pgfmathtdparse| that allows one to
parse 3d expressions. The supported vector operations are |+| (addition $+$),
|-| (subtraction $-$), |*| (multiplication of the vector by a scalar), |x|
(vector product $\times$) and |o| (scalar product). There is limited support for
ordinary \pgfname\ parsings, which are to be wrapped into |[...]|.  Admittedly,
this function is currently very poor. Enhancing it will be a top priority if
this is ever to become a CTAN library/package.

\begin{command}{\pgfmathtdparse{\marg{x}}}
   Parses 3d expressions.
\end{command}

\begin{math-function}{TDx("\mvar{vector}")}
   Yields the $x$-component of a 3d expression.
\end{math-function}

\begin{math-function}{TDy("\mvar{vector}")}
   Yields the $y$-component of a 3d expression.
\end{math-function}

\begin{math-function}{TDz("\mvar{vector}")}
   Yields the $z$-component of a 3d expression.
\end{math-function}

\begin{codeexample}[preamble={\usetikzlibrary{3dtools}}]
\begin{tikzpicture}
 \path (1,2,3) coordinate (A) (2,3,-1) coordinate (B);
 \node[align=left] {%
  $A_x=\pgfmathparse{TDx("A")}\pgfmathprintnumber\pgfmathresult$,
  $A_y=\pgfmathparse{TDy("(A)")}\pgfmathprintnumber\pgfmathresult$,
  $(A-B)_z=\pgfmathparse{TDz("(A)-(B)")}\pgfmathprintnumber\pgfmathresult$,
  $[(A)\times(B)]_x=\pgfmathparse{TDx("(A)x(B)")}\pgfmathprintnumber\pgfmathresult$};
\end{tikzpicture}
\end{codeexample}


\begin{math-function}{TDunit("\mvar{vector}")}
   Yields the result of the expression normalized to unity. If the length of the
   expression is too close to zero, a warning is issued and the unnormalized vector
   gets returned. 
\end{math-function}


\begin{math-function}{screendepth("\mvar{vector}")}
   Yields distance a coordinate is above (positive) or below (negative) the
   screen. The values are only really meaningful if the user has installed some
   reasonable view. The larger the screen depth of a point is, the closer is the
   point to the observer. The function reconstructs the 3-bein\footnote{A 3-bein
   or dreibein is a German word that stands for a local frame, its literal
   translation is something like three-leg. Like the term eigenvector this is a
   foreign word that, to the best of my knowledge, has no commonly used English
   counterpart.} from lengths like \texttt{\textbackslash pgf@xy} and so on, so the
   function is independent of the tool that is employed to install a view (cf.\
   section~\ref{sec:3DOrthonormalProjections}). The screen depth is crucial to
   decide whether a given object is in front of other objects, or behind. 
\end{math-function}

\begin{math-function}{tddistance("\mvar{point 1}","\mvar{point 2}")}
   Yields the distance between \mvar{point 1} and \mvar{point 2}. You have to
   keep the parentheses, e.g.\ |tddistance("(P)","(Q)")| works if you defined
   the points |(P)| and |(Q)|, but not 
   |tddistance("P","Q")|;															                                                                
\end{math-function}

\begin{math-function}{nscreenx}
   Yields the $x$-component of the normal on the screen. 
\end{math-function}

\begin{math-function}{nscreeny}
   Yields the $y$-component of the normal on the screen. 
\end{math-function}

\begin{math-function}{nscreenz}
   Yields the $z$-component of the normal on the screen. 
\end{math-function}

\begin{math-function}{x2d}
   Yields the $x$-component of a symbolic coordinate on the screen. 
\end{math-function}

\begin{math-function}{y2d}
   Yields the $y$-component of a symbolic coordinate on the screen. 
\end{math-function}


In order to pretty-print the result one may want to use |\pgfmathprintvector|,
and use the math function |TD| for parsing.

\begin{command}{\pgfmathprintvector\marg{x}}
   Pretty-prints vectors.
\end{command}


\begin{codeexample}[width=6.5cm,preamble={\usetikzlibrary{3dtools}}]
\pgfmathparse{TD("0.2*(A)
-0.3*(B)+0.6*(C)")}%
$0.2\,\vec A-0.3\,\vec B+0.6\,\vec C
=(\pgfmathprintvector\pgfmathresult)$
\end{codeexample}

The alert reader may wonder why this works, i.e.\ how would \tikzname\ ``know''
what the coordinates $A$, $B$ and $C$ are. It works because the coordinates in
\tikzname\ are global, so they get remembered from the above example.


\paragraph{Named coordinates.}
\label{sec:NamedCoordinates}
\tikzname\ only remembers where a named node or coordinate is on the canvas.
To know the 3d point behind it, this library records, for every named node and
coordinate, its canvas position together with its screen depth, i.e.\ its
distance from the screen (see |screendepth| below). For any non-degenerate view,
these two determine the point. The depth is tracked through
\begin{itemize}
 \item $(x,y,z)$, $(x,y)$ (which \tikzname\ reads as $(x,y,0)$) and polar
 coordinates with dimensionless radii;
 \item named coordinates and anchors: the anchors of a node have the depth of
 the node, since nodes are flat on the screen;
 \item |++(...)|, |+(...)|, |shift={(...)}| and local frames
 (section~\ref{sec:LocalFrames});
 \item sums, differences, factors and partway modifiers |!|\meta{t}|!| of the
 |calc| library;
 \item |coordinate[pos=|\meta{t}|]| on straight |--| segments.
\end{itemize}
So |\coordinate (P) at (1,2,3);|, |(A) ++(1,0,0) coordinate (P)| and
|($(A)!0.5!(B)$) coordinate (M)| all give the expected points. A coordinate in
canvas units, such as |(1cm,2cm)| or |(30:1cm)|, lies in the plane parallel to
the screen through the origin of the coordinate system in force. The depth is
\emph{not} tracked for distance and projection modifiers (|!1cm!|, |!(C)!|) and
rotations in |calc|, perpendicular coordinates |(A|\verb+|-+|B)|, |pos=| on curves,
arcs and \verb+-|+ or \verb+|-+ segments, intersections, and coordinates defined
in |canvas is xy plane at z=0| and its relatives or with |reset cm|. For such
coordinates |TD| issues a warning and falls back to the text with which the
coordinate was defined, i.e.\ to the |\tikz@dcl@coord@|\meta{coord} macro of
\tikzname, as earlier versions of this library did.

Inside the picture in which a coordinate was defined, and under the same view,
|TD("(P)")| yields $P$ in the coordinate system in force: in a scope with
|shift={(1,0,0)}| a coordinate defined outside of it appears shifted by
$(-1,0,0)$, and inside a local frame |TD| yields local coordinates. This is
what pics and path constructions that take named points need. Elsewhere, i.e.\
in node texts, outside of pictures and under a different view, |TD("(P)")|
yields $P$ in the coordinate system in which $P$ was defined. The results are
rounded to five decimals and accurate to about $10^{-5}$ ($10^{-4}$ in scaled
pictures); a coordinate that was given as a list of plain numbers is returned
exactly when |TD| is evaluated where it was defined. The |3d coordinate| key
stores the computed vector itself, and |TD| of such a coordinate always returns
that vector, which is what one wants for directions. Since the recorded position
is used, a coordinate defined with random numbers yields the same value each
time.

\begin{codeexample}[width=4.5cm,preamble={\usetikzlibrary{3dtools}}]
\begin{tikzpicture}
 \path[overlay] (0.1+rnd,0.1+rnd,0.1+rnd) 
 	coordinate (R);
 \def\ppv{\pgfmathprintvector\pgfmathresult}
 \draw[red] (R) circle[radius=0.1];	
 \node at (0,1)
  {\pgfmathparse{TD("(R)")}%
  $\vec R=(\ppv)$};	
 \draw[blue] (R) circle[radius=0.2];	
 \node at (0,0)
  {\pgfmathparse{TD("(R)")}%
  $\vec R=(\ppv)$};	
\end{tikzpicture}
\end{codeexample}

The main usage of |\pgfmathprintvector| is to strip off unnecessary zeros, which
emerge since the internal computations are largely done with the |fpu| library.

\begin{codeexample}[width=5.2cm,preamble={\usetikzlibrary{3dtools}}]
\pgfmathparse{TD("(1,0,0)x(0,1,0)")}%
$(1,0,0)^T\times(0,1,0)^T=
(\pgfmathprintvector\pgfmathresult)^T$
\end{codeexample}

\begin{codeexample}[width=5.2cm,preamble={\usetikzlibrary{3dtools}}]
\pgfmathparse{TD("(A)o(B)")}%
$\vec A\cdot \vec B=
\pgfmathprintnumber\pgfmathresult$
\end{codeexample}

Notice that, as of now, the only purpose of brackets |(...)| is to delimit
vectors. Further, the addition |+| and subtraction |-| have a \emph{higher}
precedence than vector products |x| and scalar products |o|. That is,
|(A)+(B)o(C)| gets interpreted as $(\vec A+\vec B)\cdot\vec C$, and
|(A)+(B)x(C)| as $(\vec A+\vec B)\times\vec C$.


\begin{codeexample}[width=5.2cm,preamble={\usetikzlibrary{3dtools}}]
\pgfmathparse{TD("(A)+(B)o(C)")}%
$(\vec A+\vec B)\cdot\vec C=
\pgfmathprintnumber\pgfmathresult$
\end{codeexample}

\begin{codeexample}[width=5.2cm,preamble={\usetikzlibrary{3dtools}}]
\pgfmathparse{TD("(A)+(B)x(C)")}%
$(\vec A+\vec B)\times\vec C=
(\pgfmathprintvector\pgfmathresult)$
\end{codeexample}

You can use square brackets |[...]| to separately parse subexpressions with the
standard \pgfname\ parser.

\begin{codeexample}[width=6.2cm,preamble={\usetikzlibrary{3dtools}}]
\pgfmathparse{TD("{[sin(45)]*(A)+[cos(30)]*(B)}")}%
$\sin(45)\,\vec A+\cos(30)\,\vec B=
(\pgfmathprintvector\pgfmathresult)^T$
\end{codeexample}


Moreover, any expression can only have either one |o| or one |x|, or none of
these. Expressions with more of these can be accidentally right.

\begin{math-function}{axisangles("\mvar{vector}")}
   Yields the the rotation angles that transforms the vector in the $z$-axis.
   Since an axis has a residual rotation symmetry, namely the rotation around
   this axis, only two angles are required, and thus returned. In the
   conventions of section~\ref{sec:3DOrthonormalProjections}, these are the
   angles $\phi$ and $\psi$. It corresponds to the macro
   |\tdplotgetpolarcoords|\marg{x}\marg{y}\marg{z} from the 
   \href{https://ctan.org/pkg/tikz-3dplot?lang=en}{|tikz-3dplot|} package.
\end{math-function}

\begin{codeexample}[width=5.2cm,preamble={\usetikzlibrary{3dtools}}]
\pgfmathparse{axisangles("(A)")}%
$\sphericalangle(\vec A)=
\pgfmathprintvector\pgfmathresult$
\end{codeexample}


% \tikzset{3d/projection precomputations/.code n args={3}{%
% \pgfmathsetmacro{\pgfutil@tmpa}{TD("#2-#3o#1-#3")/TD("#2-#3o#2-#3")}%
% \pgfmathsetmacro{\pgfutil@tmpb}{TD("#3+\pgfutil@tmpa*#2-\pgfutil@tmpa*#3")}%
% },3d/projection of/.style args={#1 on #2--#3}{%
% 3d/projection precomputations={#1}{#2}{#3},%
% insert path={[3d parse={(\pgfutil@tmpb)}]}
% }}
% 

% \begin{key}{/tikz/3d/projection of}
%     Allows one to compute the projection of a point on a line.
% \end{key}

Computations in 3d sometimes involve projections of a point on a line or a
plane. In 2d, projections of a point on a line can be conveniently done with the
|calc| library, but this does in general not yield the correct projection in 3d.
In order to make these projections more user friendly, |3dtools| offers the
possibility to define extended objects such as lines or planes. This is the same
concept as what is offered by plain \tikzname, where  the user can name
coordinates and nodes, and, if the |intersections| library is loaded, also
paths.

\begin{key}{/tikz/3d/plane through=\mvar{point 1} and \mvar{point 2} and \mvar{point 3}
named \mvar{name}}
    Defines a plane of name \mvar{name} by the requirement that it goes through
	the three specified points. Internally the plane is stored in terms of its
	normal and a point it goes through. 
\end{key}

\begin{key}{/tikz/3d/plane with normal=\mvar{normal} through \mvar{point 1} named \mvar{name}}
	Defines a plane of name \mvar{name} by the requirement that it goes through
	\mvar{point 1} and has the normal \mvar{normal}.  Internally the plane is
	stored in terms of this normal and point. 
\end{key}

\begin{key}{/tikz/3d/line through=\mvar{point 1} and \mvar{point 2} named \mvar{name}}
	Defines a line of name \mvar{name} by the requirement that it goes through
	\mvar{point 1} and \mvar{point 1}.   
\end{key}

\begin{key}{/tikz/3d/line with direction=\mvar{vector} through \mvar{point 1} named \mvar{name}}
	Defines a line of name \mvar{name} by the requirement that it goes through
	\mvar{point 1} and has direction \mvar{vector}.
\end{key}

\begin{key}{/tikz/3d/sphere with center=\mvar{center} and radius \mvar{radius} named \mvar{name}}
    Defines a sphere with center \mvar{center} and radius \mvar{radius}. 
\end{key}


\begin{key}{/tikz/3d/intersection of=\mvar{object 1} with \mvar{object 2}}
    Computes the intersection of a named object \mvar{object 1} with another
	object \mvar{object 2}. Currently only intersections of lines with lines and
	planes with lines are supported. There is also some support for
	intersections of lines with spheres. The intersections get also named by
	some predefined names. Note that if intersections cannot be found, warnings
	will be issued. Due to the limited precision of \TeX\ the decision whether
	or not an intersection actually exists may be wrong in close cases.
\end{key}

\begin{key}{/tikz/3d/aux keys/intersection 1  (initially I-1)}
    Predefined name for the first intersection.
\end{key}	

\begin{key}{/tikz/3d/aux keys/intersection 2  (initially I-2)}
    Predefined name for the second intersection.
\end{key}	

These extended objects can be used in projections.

\begin{key}{/tikz/3d/project=\mvar{point} on \mvar{named object}}
	Computes the projection of a point on some extended object, and inserts the
	projection in the path.   
\end{key}


The following code example illustrates the usage. It also makes use of the
|install view| key, which we describe in
section~\ref{sec:3DOrthonormalProjections}. 

\begin{codeexample}[width=4.5cm,preamble={\usetikzlibrary{3dtools}}]
\begin{tikzpicture}[3d/install view={phi=110,psi=0,theta=60},
	dot/.style={circle,inner sep=0pt,
		minimum size=2pt,fill}]
 \draw[every coordinate node/.append style={dot}]
   (2,1,2) coordinate[label=above:{$A$}] (A) -- 
   (1,2,1) coordinate[label=below:{$B$}] (B) --
   (2,0,0) coordinate[label=below:{$C$}] (C) -- cycle
   (0,0,0) coordinate[label=above:{$O$}] (O)
   (3,-2,1) coordinate[label=below:{$D$}] (D);
 \path[3d/plane through={(A) and (B) and (C) named pABC},
	3d/plane with normal={(1,1,1) through (C) named ptwo},
	3d/line through={(A) and (B) named lAB},
	3d/line through={(O) and (D) named lOD}];
 % project point on plane
 \path[3d/project={(O) on pABC}] coordinate (O');	
 % project point on line
 \path[3d/project={(C) on lAB}] coordinate (C');
 % intersection of plane and line
 \path[3d/intersection of={lOD with pABC}] coordinate (I);		  
 \draw[dashed] (C) -- (C')
  coordinate[dot,label=above right:{$C'=\pgfmathparse{TD("(C')")}%
  (\pgfmathprintvector\pgfmathresult)^T$}];
 \path (O') coordinate[dot,label=right:{$O'=\pgfmathparse{TD("(O')")}%
  (\pgfmathprintvector\pgfmathresult)^T$}];
 \path (I) coordinate[dot,label=above:{$I=\pgfmathparse{TD("(I)")}%
  (\pgfmathprintvector\pgfmathresult)^T
  $}]; 
\end{tikzpicture}
\end{codeexample}

\begin{key}{/tikz/convex hull of=\marg{list of coordinates}}
		Style that inserts a path for the 2-dimensional convex hull of a list of
		named coordinates.
\end{key}


\begin{codeexample}[width=4.5cm,preamble={\usetikzlibrary{3dtools}}]
\begin{tikzpicture}
\path foreach \X in {1,...,10}
{ (rnd*360:rnd*3) coordinate (p\X)
node[circle,inner sep=1pt,fill,label={below:$p_{\X}$}]{}};
\draw[convex hull of={p1,p2,p3,p4,p5,p6,p7,p8,p9,p10}];
\end{tikzpicture}
\end{codeexample}

\subsection{Orthonormal projections}
\label{sec:3DOrthonormalProjections}

The |3dtools| library can be used together with the
\href{https://ctan.org/pkg/tikz-3dplot?lang=en}{|tikz-3dplot|} package and/or
the |perspective| library. It also has its own means to install orthonormal projections.
Orthonormal projections emerge from subjecting 3-dimensional vectors to
orthogonal transformations and projecting them to 2 dimensions. They are not to
be confused with the perspective projections, which are more realistic and
supported by the |perspective| library. Orthonormal projections may be thought of a
limit of perspective projections at large distances, where large means that the
distance of the observer is much larger than the dimensions of the objects that
get depicted. 

\begin{key}{/tikz/3d/install view}
        Installs a 3d orthonormal projection.
\end{key}

The initial projection is such that $x$ is right an $y$ is up, as if we had no
third direction.

\begin{codeexample}[width=2cm,preamble={\usetikzlibrary{3dtools}}]
\begin{tikzpicture}[3d/install view]
 \draw[-stealth] (0,0,0) -- (1,0,0) 
  node[pos=1.2] {$x$};
 \draw[-stealth] (0,0,0) -- (0,1,0) 
  node[pos=1.2] {$y$};
 \draw[-stealth] (0,0,0) -- (0,0,1) 
  node[pos=1.2] {$z$};
\end{tikzpicture}
\end{codeexample}



The 3d-like pictures emerge by rotating the view. The conventions for the
parametrization of the orthogonal rotations in terms of three rotation angles
$\phi$, $\psi$ and $\theta$ are
\begin{equation}
 O(\phi,\psi,\theta)=\begin{pmatrix}
 c_{\phi}\,c_{\psi} & s_{\phi}\,c_{\psi} & -s_{\psi} \\
 c_{\phi}\,s_{\psi}\,s_{\theta}-s_{\phi}\,c_{\theta} & 
 s_{\phi}\,s_{\psi}\,s_{\theta}+c_{\phi}\,c_{\theta} 
 & c_{\psi}\,s_{\theta} \\
 c_{\phi}\,s_{\psi}\,c_{\theta}+s_{\phi}\,s_{\theta} 
 & s_{\phi}\,s_{\psi}\,c_{\theta}-c_{\phi}\,s_{\theta} & c_{\psi}\,c_{\theta} \\
%   c_{\phi}\,c_{\psi}
% &  s_{\psi} 
% &  -c_{\phi}\,c_{\theta}+s_{\phi}\,s_{\psi}\,s_{\theta} \\
%   -s_{\phi}\,c_{\theta}-c_{\phi}\,s_{\psi}\,s_{\theta} 
% &  c_{\psi}\,s_ {\theta}
% &  c_{\phi}\,s_{\theta}+s_{\phi}\,c_{\theta}\,s_{\psi} \\
%   -c_{\phi}\,s_{\psi}\,c_{\theta}+s_{\phi}\,s_{\theta} 
% &  c_{\psi}\,c_{\theta}
% & c_{\psi}\,c_{\theta}
\end{pmatrix}\;.\label{eq:O3}
\end{equation}
Here, $c_\phi:=\cos\phi$, $s_\phi:=\sin\phi$ and so on.
\begin{key}{/tikz/3d/phi (initially 0)}
        3d rotation angle.
\end{key}
\begin{key}{/tikz/3d/psi (initially 0)}
        3d rotation angle.
\end{key}
\begin{key}{/tikz/3d/theta (initially 0)}
        3d rotation angle.
\end{key}
The rotation angles can be used to define the view. The conventions are chosen
in such a way that they resemble those of the
\href{https://ctan.org/pkg/tikz-3dplot?lang=en}{|tikz-3dplot|} package, which
gets widely used. This matrix can be written as 
\[O(\phi,\psi,\theta)=R_x(\theta)\cdot R_y(\psi)\cdot R_z(\phi)\;,\]
where 
\begin{align*}
 R_x(\theta)&=
 \begin{pmatrix}
  1 & 0 & 0 \\
  0 & \cos (\theta ) & \sin (\theta ) \\
  0 & -\sin (\theta ) & \cos (\theta ) \\
 \end{pmatrix}\;,&
 R_y(\psi)&=
 \begin{pmatrix}
  \cos (\psi ) & 0 & -\sin (\psi ) \\
  0 & 1 & 0 \\
  \sin (\psi ) & 0 & \cos (\psi ) \\
 \end{pmatrix}\;,\\
 R_z(\phi)&=
 \begin{pmatrix}
  \cos (\phi ) & \sin (\phi ) & 0 \\
  -\sin (\phi ) & \cos (\phi ) & 0 \\
  0 & 0 & 1 \\
 \end{pmatrix}
\end{align*}
are rotations about the $x$, $y$ and $z$ axis, respectively.  For $\psi=0$,
$O(\phi=\phi_d,\psi=0,\theta=\theta_d)=R^d(\theta_d,\phi_d)$ from the
\href{https://ctan.org/pkg/tikz-3dplot?lang=en}{|tikz-3dplot|} package. Note,
however, that there seems to be an inconsistency in equation (2.1) of that
package.\footnote{I do not know how to contact the author.}

\begin{codeexample}[width=2.5cm,preamble={\usetikzlibrary{3dtools}}]
\begin{tikzpicture}[3d/install view={phi=110,psi=0,theta=70}]
 \draw[-stealth] (0,0,0) -- (1,0,0) 
  node[pos=1.2] {$x$};
 \draw[-stealth] (0,0,0) -- (0,1,0) 
  node[pos=1.2] {$y$};
 \draw[-stealth] (0,0,0) -- (0,0,1) 
  node[pos=1.2] {$z$};
\end{tikzpicture}
\end{codeexample}

\begin{key}{/tikz/3d/define orthonormal dreibein=\meta{options} (initially empty)}
		Defines a local dreibein from three input points. The initial choice of
		these points are |A=(A)|, |B=(B)|, |C=(C)|, and the initial choice for
		the basis vectors is |ex=(ex)|, |ey=(ey)|, |ez=(ez)|. In terms of these
		points, $(ex)$ is parallel to $(B)-(A)$, $(ey)$ is in the plane that
		contains $(A)$, $(B)$ and $(C)$, and is orthogonal to $(ex)$, and $(ez)$
		is orthogonal to $(ex)$ and $(ey)$. All the basis vectors are normalized
		to length 1. To draw in the plane of the dreibein, use
		|3d/install local frame=from dreibein| (section~\ref{sec:LocalFrames}).
		Using the basis vectors directly, as in |x={(ex)}|, only works if they
		have been defined without a shift because |x={(ex)}| uses the
		position of |(ex)| on the canvas.
\end{key}

\begin{key}{/tikz/3d/screen coords}
		Defines locally screen coordinate system, i.e.\ this style brings you
		back to the original coordinate system in |tikzpicture|s. It is taken
		from |tikz-3dplot|, where this style is called |tdplot_screen_coords|.
\end{key}


\subsection{Main and local frames}
\label{sec:LocalFrames}

The coordinates $(x,y,z)$ that are used after a view has been installed, e.g.\
with |3d/install view|, refer to the \emph{main} frame. Many objects are much
easier to describe in a \emph{local} (or secondary) frame. For instance, a
circle in a tilted plane is just |(0,0) circle[radius=1]| in a frame whose $xy$
plane is the tilted plane. A local frame is given by an origin $\vec O$ and
three basis vectors $\vec e_x'$, $\vec e_y'$ and $\vec e_z'$, all specified in
main coordinates. The point with local coordinates $(a,b,c)$ is
\[\vec P=\vec O+a\,\vec e_x'+b\,\vec e_y'+c\,\vec e_z'\]
in the main frame.

\begin{key}{/tikz/3d/install local frame=\meta{options}}
		Installs a local frame. It sets the |x|, |y| and |z| vectors to the
		screen images of $\vec e_x'$, $\vec e_y'$ and $\vec e_z'$, and shifts
		the origin to $\vec O$. Afterwards, |(a,b,c)| refers to the local frame,
		and so do |(a,b)|, which \tikzname\ reads as |(a,b,0)|, polar
		coordinates and circles, ellipses and arcs with dimensionless radii.
		Things that live on the canvas, such as rectangles and node texts, are
		not affected; use |canvas is xy plane at z=0| (and |transform shape|)
		inside the local frame to put them in its $xy$ plane. The vectors in
		\meta{options} are 3d expressions in the sense of |TD|, i.e.\ they can
		be explicit, like |{(1,0,0)}|, named, like |(ex)|, or combinations,
		like |(B)-(A)|. Values that contain commas must be enclosed in braces.
		The vectors always refer to the main frame, even if another local frame
		is active; local frames do not compose. Like all \tikzname\
		transformations, the key takes effect after the other options of the
		scope or path it is given to. If it fails because the basis vectors are
		(almost) coplanar, a warning is issued and nothing changes.
\end{key}

\begin{key}{/tikz/3d/local frame/origin=\meta{point} (initially {(0,0,0)})}
		Origin of the local frame.
\end{key}
\begin{key}{/tikz/3d/local frame/ex=\meta{vector} (initially {(1,0,0)})}
		Local $x$ basis vector.
\end{key}
\begin{key}{/tikz/3d/local frame/ey=\meta{vector} (initially {(0,1,0)})}
		Local $y$ basis vector.
\end{key}
\begin{key}{/tikz/3d/local frame/ez=\meta{vector} (initially empty)}
		Local $z$ basis vector. If empty, $\vec e_z'=\vec e_x'\times\vec e_y'$.
\end{key}
\begin{key}{/tikz/3d/local frame/orthonormalize=\meta{boolean} (default true, initially false)}
		If true, the basis is replaced by the right-handed orthonormal basis
		with $\vec e_x'\parallel$ |ex| and $\vec e_y'$ in the plane of |ex|
		and |ey|; |ez| is ignored. Otherwise the vectors are used as they are,
		so the local frame may be skewed or scaled.
\end{key}
\begin{key}{/tikz/3d/local frame/name=\meta{name} (initially empty)}
		Saves the frame under \meta{name} for later use with |3d/use local frame|.
\end{key}
\begin{key}{/tikz/3d/local frame/from dreibein}
		Uses the frame defined by the last |3d/define orthonormal dreibein|,
		i.e.\ origin |A| and basis vectors |ex|, |ey| and |ez|.
\end{key}

\begin{key}{/tikz/3d/use local frame=\meta{name}}
		Installs a local frame that has been saved with |name=|\meta{name}.
\end{key}

\begin{key}{/tikz/3d/main frame}
		Goes back to the main frame inside a local frame, e.g.\ for a single
		path. The main frame consists of the |x|, |y| and |z| vectors and the
		transformation that were in force when the first local frame of the
		current scope got installed.
\end{key}

\begin{codeexample}[width=5cm,preamble={\usetikzlibrary{3dtools}}]
\begin{tikzpicture}[3d/install view={phi=110,psi=0,theta=70},
  dot/.style={circle,inner sep=1pt,fill}]
 \draw[-stealth] (0,0,0) -- (3,0,0) node[pos=1.1]{$x$};
 \draw[-stealth] (0,0,0) -- (0,3,0) node[pos=1.1]{$y$};
 \draw[-stealth] (0,0,0) -- (0,0,3) node[pos=1.1]{$z$};
 \path (2,0,0) coordinate (A) (0,2,0) coordinate (B)
  (0,0,2) coordinate (C) (2/3,2/3,2/3) coordinate (I);
 \draw[gray] (A) -- (B) -- (C) -- cycle;
 \begin{scope}[3d/install local frame={origin=(I),
   ex=(B)-(A),ey=(C)-(A),orthonormalize,name=tilted}]
  \draw[blue] (0,0) circle[radius=0.8];
  \draw[blue] (-.5,-.5) -- (.5,-.5) -- (.5,.5)
   -- (-.5,.5) -- cycle;
  \draw[-stealth] (0,0,0) -- (0,0,1) node[above]{$\vec n$};
  \begin{scope}[canvas is xy plane at z=0,transform shape]
   \node[blue] at (0,-0.6) {\tiny tilted};
  \end{scope}
  \path (0.8,0) coordinate[dot,label=right:$P$] (P);
 \end{scope}
\end{tikzpicture}
\end{codeexample}

Inside the picture, |TD("(P)")| yields the main coordinates of $P$ outside
the local frame and its local coordinates $(0.8,0,0)$ inside it
(section~\ref{sec:NamedCoordinates}). Outside the picture, as in the next
example, it yields the coordinates $P$ was defined with, i.e.\ the local ones.
The following functions convert explicit coordinates between the frames. They use the local frame that is
in force or, outside of any local frame, the one that was installed last.
Results are accurate to about $10^{-4}$.

\begin{math-function}{TDtomain("\mvar{vector}")}
   Yields the main coordinates of a point given in local coordinates.
\end{math-function}

\begin{math-function}{TDtolocal("\mvar{vector}")}
   Yields the local coordinates of a point given in main coordinates.
\end{math-function}

\begin{codeexample}[width=5.6cm,preamble={\usetikzlibrary{3dtools}}]
\pgfmathsetmacro{\Pmain}{TDtomain("(P)")}%
\pgfmathsetmacro{\Iloc}{TDtolocal("(I)")}%
\pgfmathsetmacro{\Aloc}{TDtolocal("(A)")}%
\def\ppv#1{(\pgfmathprintvector[/pgf/number format/fixed]{#1})}%
$\begin{array}{l}
\vec P_\mathrm{main}=\ppv{\Pmain},\\
\vec I_\mathrm{loc}=\ppv{\Iloc},\\
\vec A_\mathrm{loc}=\ppv{\Aloc}
\end{array}$
\end{codeexample}

Pics work in local frames, too. Here a flag pic, which is drawn in its own $xy$
plane, gets placed at six positions on a circle.

\begin{codeexample}[width=5cm,preamble={\usetikzlibrary{3dtools}}]
\begin{tikzpicture}[3d/install view={phi=120,psi=0,theta=65},
 pics/flag/.style={code={\draw[thick] (0,0) -- (0,1.3);
   \draw[fill=yellow!60] (0,0.7) -- (0.8,0.7) -- (0.8,1.25)
    -- (0,1.25) -- cycle;}}]
 \draw (0,0,0) circle[radius=1.5];
 \foreach \a in {0,60,...,300}
 {\begin{scope}[3d/install local frame={
    origin={({1.5*cos(\a)},{1.5*sin(\a)},0)},
    ex={({cos(\a)},{sin(\a)},0)},ey={(0,0,1)}}]
   \pic{flag};
  \end{scope}}
\end{tikzpicture}
\end{codeexample}

\subsection{Visibility considerations}

When drawing objects in 3d, some parts are ``visible'', i.e.\ in the foreground,
while other parts are ``hidden'', i.e.\ in the background. In what follows, we
discuss some basic means that allow one to determine whether some stretch is
visible or hidden. Very often this amounts to find critical values of some
parameter(s) (such as angles) at which some curve changes from visible to
hidden, or vice versa. One of the major players in this game is the normal on
the screen, which is given by the last row of the orthogonal matrix $O$ of
\eqref{eq:O3}. For $\psi=0$, this normal is also the normal in the so-called
``main'' coordinates of 
\href{https://ctan.org/pkg/tikz-3dplot?lang=en}{|tikz-3dplot|} with the
identifications $\theta=\texttt{\textbackslash tdplotmaintheta}$ and
$\phi=\texttt{\textbackslash tdplotmainphi}$.

\begin{key}{/tikz/3d/draw ordered paths=\meta{list} (initially empty)}
		This draws a list of named paths that have been created in with the
		|save named path| in the order they appear in the list \meta{list}.
		This is achieved by a number of clips and inverted clips. Notice that,
		at this point, layers are not supported. This is because clips have to
		be applied in the some layer as the path that is to be clipped against.
		One can define styles that carry the same names as the named paths in
		the directory |3d/ordered paths/|. If such a style exists, it will be
		applied. 
\end{key}

\begin{codeexample}[width=5cm,preamble={\usetikzlibrary{3dtools}}]
\begin{tikzpicture}[3d/install view={phi=50,theta=70},
  declare function={a=3;alpha=20;},
  line join=round]
 \path
   (-a/2,0,0) coordinate (P_1)
   (a/2,0,0) coordinate (P_2)
   (-a/2,a,0) coordinate (A_1)
   (a/2,a,0) coordinate (A_2)
   (-a/2,{a*cos(alpha/2)},{a*sin(alpha/2)}) coordinate (B_1)
   (a/2,{a*cos(alpha/2)},{a*sin(alpha/2)}) coordinate (B_2)
   (-a/2,{a*cos(alpha)},{a*sin(alpha)}) coordinate (C_1)
   (a/2,{a*cos(alpha)},{a*sin(alpha)}) coordinate (C_2);
 \foreach \X in {A,B,C}	 	 
 {\path[save named path=p\X] (P_1) -- (P_2) -- 
 	(\X_2) -- (\X_1) -- cycle;}	 
 \pgfmathsetmacro{\nA}{TDunit("(A_1)-(P_1)x(P_2)-(P_1)")}
 \pgfmathsetmacro{\nC}{TDunit("(C_1)-(P_1)x(P_1)-(P_2)")}
 \pgfmathtruncatemacro{\itest}{screendepth(\nA)>screendepth(\nC)?1:0}
 \tikzset{3d/ordered paths/pA/.style={fill=red}}
 \ifnum\itest=1
  \tikzset{3d/draw ordered paths={pC,pB,pA}}
 \else
  \tikzset{3d/draw ordered paths={pA,pB,pC}}
 \fi
\end{tikzpicture}
\end{codeexample}

\subsubsection{Example 1: latitude circle on sphere}

Here is an example. Suppose we wish to draw a latitude circle on a sphere for
$\psi=0$. Since the system has a rotational symmetry around the $z$ axis, we can
set $\phi=0$ as well. The normal on the screen then becomes
\[
 \vec n_\mathrm{screen}=
 \bigl(0,-\sin(\theta),\cos(\theta)\bigr)^T\;.
\]
The latitude circle can be parametrized as
\[
 \vec\gamma_\mathrm{lat}(\beta)=
 \bigl(r\,\cos(\beta),r\,\sin(\beta),h\bigr)^T\;,
\]
where $r=R\,\cos(\lambda)$ and $h=R\,\sin(\lambda)$ with $R$ being the radius of
the sphere and $\lambda$ the latitude angle of the latitude circle. The critical
angles $\beta_\mathrm{crit}$ at which the curve transitions from hidden to
visible are then given by the solutions of
\[
 \vec n_\mathrm{screen}\cdot\vec\gamma_\mathrm{lat}(\beta_\mathrm{crit})
 =-r\,\sin(\beta)\,\sin(\theta)+h\,\cos(\theta)
 \overset{!}{=}0\;.
\]
Since $-1\le\sin(\beta)\le1$, this equation only has solutions if 
\begin{equation}\label{eq:disclat}
 \left\vert\frac{h}{r}\cot(\theta)\right\vert\le1\;.
\end{equation} 
The solutions are then given by 
\[
 \beta_\mathrm{crit}^{(1)}=\arcsin\left(\frac{h}{r}\cot(\theta)\right)
 \quad\text{and}\quad
 \beta_\mathrm{crit}^{(2)}=180-\beta_\mathrm{crit}^{(1)}\;,
\]
and $\beta_\mathrm{crit}^{(1)}=\beta_\mathrm{crit}^{(2)}$ if the inequality 
\eqref{eq:disclat} is saturated. The following code example illustrates this.
Notice that it could be made more concise with the |3d| library.

\begin{codeexample}[width=5.6cm,preamble={\usetikzlibrary{3dtools}}]
\begin{tikzpicture}[declare function={%
	R=2.5;lambda=20;theta=60;
	r=R*cos(lambda);h=R*sin(lambda);
	betacrit=asin(h/r*cot(theta));}]
 \draw (0,0) circle[radius=R];	
 \begin{scope}[smooth,
 	3d/install view={phi=0,psi=0,theta=theta}]
  \draw[dashed] 
   plot[variable=\t,
   	domain=betacrit:180-betacrit]
   ({r*cos(\t)},{r*sin(\t)},h);
  \draw 
   plot[variable=\t,
   	domain=betacrit:-180-betacrit]
   ({r*cos(\t)},{r*sin(\t)},h);
 \end{scope}
\end{tikzpicture}
\end{codeexample}


\subsubsection{Example 2: arbitrary circle on sphere}

A somewhat more advanced question is to draw an arbitrary circle on a sphere.
That is, we are given a sphere around a center $C$ or radius $R$ and a point $P$
inside the sphere. If $C$ and $P$ coincide, we need to define a normal $n$,
otherwise $n=P-C$.  How can one draw a circle on the sphere, distinguishing
between visible and hidden stretches?

\paragraph{\boldmath First test: $\lvert P-C\rvert<R$.}
 The circle only makes sense if
$\vert P-C\vert<R$. The radius of the circle is $r=\sqrt{R^2-\lvert P-C\rvert^2}$. Call the normal
on the screen $n_\mathrm{screen}$. 

\paragraph{\boldmath Case $n\parallel n_\mathrm{screen}$.}  If $n$ and
$n_\mathrm{screen}$ are linearly dependent, the circle is completely hidden if
the screen depth of $P$ is smaller than the screen depth of $C$, $\sd P <\sd$,
and visible otherwise. We can choose a coordinate system in which  
\[
e_z=\frac{n}{|n|}\;,\quad e_x\text{ arbitrary but }\perp e_z \quad\text{and}\quad
e_y=e_z\times e_x
\] 
to draw the circle.

\paragraph{\boldmath Case $n\not\,\parallel n_\mathrm{screen}$.} 
From now on we assume that $n$ and $n_\mathrm{screen}$ are not linearly
dependent. This means that they span a plane, and we can choose
$n_\mathrm{screen}$ to point in the $x$-direction of this plane (see
figure~\ref{fig:not_parallel}).

\begin{figure}[htb]
\centering
\begin{tikzpicture}
 \draw (-3,-3) rectangle (4,3);
 \draw[dashed] (0,3) coordinate (t) -- (0,-3) coordinate (b)node[below]{$\sd C$};
 \draw[-stealth] (3.5,-3.3) -- (4,-3.3) node[right] {$\sd$};
 \path (2,-2) coordinate (c1) -- (-1,2) coordinate (c2)
 	node[midway,circle,fill,inner sep=1.2pt,label=above:{$P$}](P){};
 \draw[thick,blue] (intersection of c1--c2 and t--b) coordinate (i) -- coordinate (v) (c1) ;	
 \draw[thick,gray] (i) -- coordinate (h) (c2) ;	
 \draw ($(c1)!5mm!90:(c2)$) --  coordinate (r1) ($(c1)!2mm!90:(c2)$)
 ($(c2)!5mm!-90:(c1)$) --  coordinate (r2) ($(c2)!2mm!-90:(c1)$)
 (r1) -- node[sloped,fill=white]{$r$} (r2);
 \draw[red,-stealth] (P) -- ($(P)!20mm!90:(c1)$)coordinate (n) node[pos=1.15]{$n$};
 \draw[red,dashed] (n) |- (P) node[pos=0.25,right]{$n_\perp$}
  coordinate[pos=0.5] (aux)
  node[pos=0.75,above]{$n_\parallel$}
  pic[draw,solid,pic text={$\alpha$},angle radius=3em,angle eccentricity=0.6]{angle=c1--P--aux};
\end{tikzpicture}
\caption{$n\not\,\parallel n_\mathrm{screen}$.}
\label{fig:not_parallel}
\end{figure}

The normal can be decomposed in a part parallel to $n_\mathrm{screen}$,
$n_\parallel$, and a perpendicular part, $n_\perp$. The projection of the circle
on this plane is a line, and has slope $\alpha=\arctan(-n_\parallel/n_\perp)$.
If $r\cos\alpha>\lvert\sd P-\sd C\lvert$, the circle has both hidden and visible
stretches, from between $-r\cos\alpha$ and $\sd P-\sd C$ it is hidden and between $\sd
P-\sd C$ and $r\cos\alpha$ it is solid. 

We are going to draw the circle in a coordinate system defined by
\[ e_z=\frac{n}{\lvert n \rvert}\;,\quad e_y=e_z\times n_\mathrm{screen}
\quad\text{and}\quad e_x=e_y\times e_z\;.\]
The circle is then in the $x$-$y$ plane, and the positive $x$-direction is the
direction of increasing screen depth (or visibility). If $r\cos\alpha>\lvert \sd P\rvert$,
then the stretch between $\beta:=\left\vert\arccos\left(\frac{\sd C-\sd
P}{r\cos\alpha}\right)\right\vert$ and $-\beta$ is visible, between $\beta$ and
$360^\circ-\beta$ it is hidden.

\subsubsection{Example 3: outlines of shapes}

Another type of visibility considerations concerns the contours of shapes. To be
specific, consider a surface parametrized by
\begin{equation}\label{eq:f}
 \vec f(u,v)=\begin{pmatrix}
 u\\ R(u)\,\cos (v)\\ R(u)\,\sin (v)
 \end{pmatrix}
\end{equation}
with some function $R(u)$. What is the contour of an orthonormal projection of the
surface on the screen? It is the boundary between the hidden and visible parts.
What distinguishes the visible from the hidden parts? It is the projection of
the normal of the surface on the screen, $F(u,v):=n_S(u,v)\cdot 
n_\mathrm{screen}$. The hidden and visible stretches are separated by the zeros
of this projection. In the case of our surface \eqref{eq:f},
\begin{equation}
  \vec n_S(u,v)=\vec\nabla_u f(u,v)\times \vec\nabla_v f(u,v)
 =R(u)\,\begin{pmatrix}
 	-R'(u)\\ \cos (v)\\ \sin (v)
 \end{pmatrix}\;,
\end{equation}
and for $\psi=0$
\begin{equation}
 \vec n_\mathrm{screen}=\begin{pmatrix}
 	\sin(\theta)\,\sin (\phi)\\
	-\sin (\theta )\,\cos(\phi)\\
	\cos (\theta)
  \end{pmatrix}\;.
\end{equation}
So we need to find the zeros of 
\begin{align}
 F(u,v)={}&{}-\cos (\theta)\,\sin (v)+
 \sin (\theta )\,\cos(\phi)\,\cos (v)+ R'(u)\,\sin(\theta)\,\sin (\phi)\notag\\
 =:{}&{}a\,\sin (v)+b\,\cos(v)+c\;.\label{eq:vcrit}
\end{align} 
This equation can be solved for $x:=\sin (v)$ using 
$\cos(v)=\pm\sqrt{1-\sin^2(v)}=\pm\sqrt{1-x^2}$. It is clear that, since  $\sin
(v)$ is bounded (and so is $\cos(v)$), this equation does not always have a
solution. In particular, if $R'(u)$ is too large, there might not be a solution.
This just means that for a given $u$ all points are hidden or visible. This
happens e.g.\ when you look at a cone from the top or bottom. A more advanced
code could deal with these cases, the one presented below (which is taken from
\href{https://topanswers.xyz/tex?q=1218\#a1447}{topanswers.xyz}) does not. Note also
that these are local visibility conditions, a seemingly visible stretch can
always be covered by some other stretch that is a finite distance away in $u$
direction.

However, for ``reasonable'' configurations, there are two solutions of the
quadratic equation, as they should, they correspond to the upper and lower
contour in the code example below. They are a bit unilluminating and given by
\begin{subequations}\label{eq:cvandsv}
\begin{align}
 [\sin(v)]_{1,2}&=\frac{-a\,c\pm b \sqrt{a^2+b^2-c^2}}{a^2+b^2}\;,\\
 [\cos(v)]_{1,2}&=\frac{-b\,c\mp a \sqrt{a^2+b^2-c^2}}{a^2+b^2}\;,
\end{align}
\end{subequations}
where $a$, $b$ and $c$ are implicitly defined in \eqref{eq:vcrit}. Hence, for
``reasonable'' configurations, it is straightforward to plot the contour of the
surface with \LaTeX.


\begin{codeexample}[width=2.5cm,preamble={\usetikzlibrary{3dtools}}]
\begin{tikzpicture}[declare function={phi=30;theta=70;},
	3d/install view={phi=phi,psi=0,theta=theta},
	line cap=butt,line join=round,
	declare function={%
	  R(\u)=1.5+sin(\u*45)/(2+\u/8);%<- input
	  Rprime(\u)=(R(\u+0.01)-R(\u-0.01))/0.02;%<- numerical derivative
	  a=cos(phi)*sin(theta);% see equation {eq:vcrit}
	  b=-1*cos(theta);%	
	  c(\u)=sin(theta)*sin(phi)*Rprime(\u);%
      sv1(\u)=-((b*c(\u)+sqrt(a*a+b*b-c(\u)*c(\u))*abs(a))/(a*a+b*b));% 
	  sv2(\u)=(-(b*c(\u))+sqrt(a*a+b*b-c(\u)*c(\u))*abs(a))/(a*a+b*b);%
	  cv1(\u)=-((a*c(\u)+sqrt(a*a+b*b-c(\u)*c(\u))*abs(b))/(a*a+b*b));%
	  cv2(\u)=(-(a*c(\u))+sqrt(a*a+b*b-c(\u)*c(\u))*abs(b))/(a*a+b*b);%
	},
	icircle/.style={thick},
	iplane/.style={fill=white,fill opacity=0.8},
	iline/.style={semithick},
 	extra/.code={},annot/.code={\tikzset{extra/.code={#1}}}]
 \newcommand\DrawIntersectingPlane[2][]{%
  \begin{scope}[canvas is yz plane at x=#2,transform shape,#1]
	 \draw[iplane] (3,3) -- (-3,3) -- (-3,-3)-- (3,-3) -- cycle;
	 \draw[icircle]  (0,0) circle[radius={R(#2)}];
	 \tikzset{extra}
  \end{scope}}
 %
 \DrawIntersectingPlane[annot={%
 	\path (2.9,2.9) node[below left,transform shape]{$P$};
	\draw (2,3) arc[start angle=180,end angle=270,radius=1];
	\draw[dashed] (0,-2.25) node[below,transform shape=false]{$a$} -- (0,0);
	}]{0}
 \draw[iline] (0,0,-2.25) -- (4,0,-2.25);
 \draw[thick,smooth,variable=\u,domain=0:4] 
 	plot  (\u,{R(\u)*cv1(\u)},{R(\u)*sv1(\u)})
 	plot  (\u,{R(\u)*cv2(\u)},{R(\u)*sv2(\u)});
 \DrawIntersectingPlane[icircle/.append style={fill=gray,fill opacity=0.5},
   annot={%
	\draw[dashed] (0,-2.25) node[below,transform shape=false]{$x$} -- (0,-1.5);
	}]{4}
 \draw[iline] (4,0,-2.25) -- (8,0,-2.25);
 \draw[thick,smooth,variable=\u,domain=4:8] 
 	plot  (\u,{R(\u)*cv1(\u)},{R(\u)*sv1(\u)})
 	plot  (\u,{R(\u)*cv2(\u)},{R(\u)*sv2(\u)});
 %
 \DrawIntersectingPlane[annot={%
 	\path (2.9,2.9) node[below left,transform shape]{$Q$};
	\draw (2,3) arc[start angle=180,end angle=270,radius=1];
	\draw[dashed] (0,-2.25) node[below,transform shape=false]{$b$} -- (0,0);
	}]{8}
 \draw[iline] (8,0,-2.25) -- (9,0,-2.25) ;
\end{tikzpicture}
\end{codeexample}
% \caption{Example of a screen projection of a surface of revolution. 
% Taken from \href{https://topanswers.xyz/tex?q=1218\#a1447}{topanswers.xyz}.}
% \label{fig:example}
% \end{figure}

There are some special cases that may be of particular interest: $R'(u)=0$ for a
cylinder and $R'(u)=-r/h$ for a cone of height $h$ and base radius $r$. 

\subsection{Automatic occlusion}
\label{sec:Scenes}

TikZ paints in the order in which paths are given. Inside a scene, the
library reorders them instead: every path becomes a \emph{piece}, its output is
put aside while TikZ builds it, and at the end of the scene the pieces are
painted from back to front (painter's algorithm). The depth of a piece is the
mean depth of the points the path visits (section~\ref{sec:NamedCoordinates}).
The pics |cylinder|, |3d/frustum|, |3d/cone|, |3d/shaded cone|, |3d/circle|
and |3d/circle on sphere| split themselves into pieces (the rim behind, the
front half of the mantle, the cap, the arcs in front and behind), each with the
depth of its centroid, so that, e.g., a ring around a cylinder passes behind
it and in front of it. Faces of polyhedra are paths, so they are pieces, too.

A straight path whose depth varies (an axis, a beam) is cut where it passes
other pieces (at their depths), and each run is painted at its own depth, so it
can pass through them (key |3d/split|). Pieces of equal depth, like the
coplanar caps of nested cylinders, are painted larger first. A |3d/sphere|
places the pieces around it: what is inside stays between its back and front
shells, what lies on its surface (points, the arcs of
|3d/circle on sphere|) goes just in front of the shell it lies on, in the
order of the code on the front and in the reverse order on the back (what is
on top in front is at the bottom behind), and what is outside goes behind or
in front of the whole sphere. A straight path is also cut where it passes the
surface of a sphere. In a scene the hidden edges
of the pics are drawn like visible ones: an opaque surface covers them, a
translucent one tints them.

This is a heuristic, not hidden surface removal. Each piece gets a single
depth, so curved pieces that interpenetrate may get the wrong order; set the
depth by hand with |3d/depth| where needed. A circle drawn with the |circle|
path operation is a single piece at the depth of its center; use the pic
|3d/circle| for rings that wrap around other objects, and |3d/sphere| for a
sphere with things inside. Translucent fills keep working.

\begin{key}{/tikz/3d/scene}
	Turns the current scope (or picture) into a scene. Nested scenes belong to
	the outermost one. A piece goes to the layer that was active when the path
	started. The line width, dash pattern, colors, opacity, line cap and join
	set by scopes inside the scene are restored in every piece; clips and
	transparency groups of such scopes are not, and should be put on the paths
	(or around the scene). Paths that visit no point with a known depth, like a
	lone node at a named coordinate, are painted in front.
\end{key}

\begin{key}{/tikz/3d/split=\meta{length} (initially 0.2mm)}
	A path whose points lie on a straight line, that does not fill, and whose
	depth varies by more than \meta{length} is cut where it passes other pieces
	(the runs are clipped copies of the same drawing, overlapping by 0.25pt).
	|0| switches this off.
\end{key}

\begin{key}{/tikz/3d/edge only}
	Removes the fill, the shading and the path picture of a path. The pics
	add it to |3d/hidden|, so the style of a surface can be reused for its
	hidden edges.
\end{key}

\begin{key}{/tikz/3d/depth=\meta{depth}}
	Sets the depth of the pieces of the current path or scope. \meta{depth} is
	|front|, |back|, or a coordinate in parentheses whose depth is used, e.g.\
	|3d/depth={(0,0,1)}|.
\end{key}

\begin{codeexample}[width=6cm,preamble={\usetikzlibrary{3dtools}}]
\begin{tikzpicture}[3d/install view={phi=20,theta=70},
  3d/scene]
 % the rings come first in the code
 \foreach \z in {0.6,1.8}
  \pic[orange,ultra thick] at (0,0,\z) {3d/circle={r=1}};
 \pic {cylinder={r=0.5,h=2.4,
  cylinder/mantle/.style={fill=cyan!20,fill opacity=0.5,draw},
  cylinder/top/.style={fill=cyan!20,fill opacity=0.5,draw}}};
 \pic at (2.5,0,0) {3d/frustum={R=1,r=0.3,h=1.5,
  frustum/inner/.style={fill=yellow!70,fill opacity=0.7},
  frustum/outer/.style={fill=yellow!40,fill opacity=0.4}}};
 \draw[red,thick,->] (2.5,0,-0.5) -- (2.5,0,2.2);
\end{tikzpicture}
\end{codeexample}

\subsection{``Physical'' coordinates}
\label{sec:PhysicalCoordinates}

It is quite conceivable that users may want to use different coordinate systems
for defining different coordinates. The following functions aim at supporting
this. It should be said, though, that these functions are even more
``experimental'' than what has been discussed thus far. The main point is that,
apart from the ``physical'' $x$ and $y$ coordinates, i.e.\ the coordinates of an
point on the screen, one can also keep track of the physical $z$ coordinate, or
screen depth, which indicates how far a given point is above  ($z>0$) or below
($z<0$) the screen. The physical components of every named node and coordinate
are recorded (section~\ref{sec:NamedCoordinates}).

\begin{key}{/tikz/3d/record physical components}
        Does nothing; kept for compatibility. Earlier versions needed it to
        switch on the recording.
\end{key}

The physical components of a point are its $x$ and $y$ coordinates on the
canvas, in |cm|, and its screen depth $\vec n\cdot\vec P$, where $\vec P$ are the
main coordinates of the point (section~\ref{sec:LocalFrames}) and $\vec n$ is the
normal on the screen of the main frame (cf.\ |nscreenx| etc.). For views
installed with |3d/install view| this means that the physical components are
$O(\phi,\psi,\theta)\cdot\vec P$ of \eqref{eq:O3}, plus the canvas shift of the
main frame. The rules for the depth are those of
section~\ref{sec:NamedCoordinates}; if the depth of a coordinate is not known,
|TDphysz| warns and yields 0. The physical components can then be retrieved with various functions.

\begin{math-function}{TDphysx("\mvar{vector}")}
   Yields the physical $x$-component of a 3d named coordinate.
\end{math-function}

\begin{math-function}{TDphysy("\mvar{vector}")}
   Yields the physical $y$-component of a 3d named coordinate.
\end{math-function}

\begin{math-function}{TDphysz("\mvar{vector}")}
   Yields the physical $z$-component of a 3d named coordinate.
\end{math-function}

\begin{math-function}{TDphys("\mvar{vector}")}
   Yields an array containing physical $x$, $y$ and $z$-components of a 3d named
   coordinate.
\end{math-function}

Sometimes one may want to retrieve the components of a physical coordinate in
another frame, i.e.\ the coordinates of a point that has been defined in one
frame with respect to the frame that is in force now. These functions invert
the map from coordinates to physical components exactly, so they work for any
non-degenerate view and local frame. Shifts are taken into account; other
transformations, e.g.\ |scale|, should be the same as when the point was
recorded.

\begin{math-function}{TDlocalx("\mvar{vector}")}
   Yields the local $x$-component of a 3d named coordinate the physical
   coordinates of which have been recorded.
\end{math-function}

\begin{math-function}{TDlocaly("\mvar{vector}")}
   Yields the local $y$-component of a 3d named coordinate the physical
   coordinates of which have been recorded.
\end{math-function}

\begin{math-function}{TDlocalz("\mvar{vector}")}
   Yields the local $z$-component of a 3d named coordinate the physical
   coordinates of which have been recorded.
\end{math-function}

\begin{math-function}{TDlocal("\mvar{vector}")}
   Yields an array containing local $x$, $y$ and $z$-components of a 3d named
   coordinate the physical coordinates of which have been recorded.
\end{math-function}

In the following example, $B$, $C$ and $D$ are defined at $(3,0,0)$ in views
that are rotated by $\phi=30^\circ$, $150^\circ$ and $270^\circ$, respectively.
|TDlocal|, evaluated in the unrotated view, yields their coordinates in this view.

\begin{codeexample}[width=5.2cm,preamble={\usetikzlibrary{3dtools}}]
\begin{tikzpicture}[3d/record physical components,
	dot/.style={circle,inner sep=1pt,fill}]
 \begin{scope}[3d/install view={phi=0,psi=0,theta=70}]
  \path (0,0,3) coordinate (A);
  \path[3d/install view={phi=30}] (3,0,0) coordinate (B);
  \path[3d/install view={phi=150}] (3,0,0) coordinate (C);
  \path[3d/install view={phi=270}] (3,0,0) coordinate (D);
  \draw[dashed] foreach \X [remember=\X as \Y (initially D)]
 	in {B,C,D}
   {(A) -- (\X) (\Y) -- (\X)
  coordinate[dot,label=below:{$\begin{aligned}
	\X_\mathrm{phys}&=\pgfmathparse{TDphys("\X")}%
  (\pgfmathprintvector[/pgf/number format/fixed]\pgfmathresult)^T\\
	\X_\mathrm{loc}&=\pgfmathparse{TDlocal("\X")}%
  (\pgfmathprintvector[/pgf/number format/fixed]\pgfmathresult)^T\\
  \end{aligned}$}]}
  (A) coordinate[dot,label=above:{$A=\pgfmathparse{TDphys("A")}%
  (\pgfmathprintvector\pgfmathresult)^T
  $}];
 \end{scope}
\end{tikzpicture}
\end{codeexample}

Once one has defined physical coordinates, one can use them to compute e.g.\
distances between coordinates defined in different frames.
\begin{codeexample}[width=4.6cm,preamble={\usetikzlibrary{3dtools}}]
$\pgfmathsetmacro{\vecA}{TDphys("A")}%
\pgfmathsetmacro{\vecB}{TDphys("B")}%
\pgfmathsetmacro{\physdist}{%
sqrt(TD("(\vecA)-(\vecB)o(\vecA)-(\vecB)"))}%
\begin{array}{l}
\vec A_\mathrm{phys}=(\pgfmathprintvector\vecA),\\
\vec B_\mathrm{phys}=(\pgfmathprintvector\vecB),\\
d_\mathrm{phys}(\vec A,\vec B)=%
\pgfmathprintnumber\physdist
\end{array}$
\end{codeexample}

\begin{key}{/tikz/3d/install dual basis}
		Defines the coordinates |(dualex)|, |(dualey)| and |(dualez)| (the names
		are stored in the keys |/tikz/3d/dual ex| etc.) such that the local
		components are the scalar products of these vectors with the physical
		components measured from the origin of the frame in force. For
		orthonormal views and frames, they are the physical components of the
		basis vectors.
\end{key}

Section~\ref{sec:NamedCoordinates} lists the constructions through which the
depth cannot be tracked, e.g.\ |canvas is xy plane at z=0| (use a local frame
instead).

\subsection{Predefined path constructions and pics}

\begin{key}{/tikz/3d/circumsphere center=\meta{options} (initially empty)}
		Computes the center of a sphere for a given set of four non-coplanar points. The
		initial choice of these points are |A=(A)|, |B=(B)|, |C=(C)|, |D=(D)|.
		The underlying maths can be found at
		\href{https://topanswers.xyz/tex?q=1233#a1467}{https://topanswers.xyz/tex?q=1233\#a1467}.
\end{key}

\begin{codeexample}[width=2.5cm,preamble={\usetikzlibrary{3dtools}}]
\begin{tikzpicture}[dot/.style={circle,inner sep=1pt,fill}]
 \begin{scope}[3d/install view={phi=100,psi=0,theta=70}]	
  \path (1,0,0) coordinate (X)
   (0,1,0) coordinate (Y)
   (0,1,1) coordinate (Z)
   (1,-1,1) coordinate (P);
  \path[3d/circumsphere center={A={(X)},B={(Y)},C={(Z)},D={(P)}}] 
  coordinate (I); 
 \end{scope}
  \pgfmathsetmacro{\csr}{sqrt(TD("(X)-(I)o(X)-(I)"))}
  \draw (I) circle[radius=\csr];  
  \path foreach \X in {X,Y,Z,P}
 {(\X) coordinate[dot,label=below:{$\X$}]}
 (I) (I) coordinate[dot,label=above:{$I=\pgfmathparse{TD("(I)")}%
  (\pgfmathprintvector\pgfmathresult)^T$}];
\end{tikzpicture}
\end{codeexample}

\begin{key}{/tikz/3d/circumcircle center=\meta{options} (initially empty)}
		Computes the center of a circle for a given set of three non-collinear
		points. The initial choice of these points are |A=(A)|, |B=(B)|,
		|C=(C)|.
\end{key}

\begin{codeexample}[width=5.4cm,preamble={\usetikzlibrary{3dtools}}]
\begin{tikzpicture}[declare function={a=3;},
   dot/.style={circle,inner sep=1pt,fill},
   3d/install view={phi=100,psi=0,theta=70}]	
  \path (a,0,0) coordinate (A)
   (0,a,0) coordinate (B)
   (0,0,a) coordinate (C);
  \path[3d/circumcircle center] 
  	coordinate (I); 
  \pgfmathsetmacro{\myr}{%
     sqrt(TD("(A)-(I)o(A)-(I)"))}
  \tikzset{3d/define orthonormal dreibein}
  \begin{scope}[x={(ex)},y={(ey)}]
   \draw (I) circle[radius=\myr];
  \end{scope}
  \path foreach \X in {A,B,C}
  {(\X) coordinate[dot,label=above:{$\X$}]}
  (I) (I) coordinate[dot,label=above:{$I=%
  	\pgfmathparse{TD("(I)")}%
   (\pgfmathprintvector\pgfmathresult)^T$}];
\end{tikzpicture}
\end{codeexample}

\begin{key}{/tikz/3d/intersection of three spheres=\meta{options} (initially empty)}
		Computes the intersection intersections of three spheres around
		|A=(A)|, |B=(B)|, and |C=(C)|. The radii are stored in |rA|, |rB| and
		|rC|, respectively. The names of the intersection solutions are stored
		in the keys |i1| and |i2|, the initial values of which are |i1| and
		|i2|.
\end{key}

\begin{codeexample}[width=5.4cm,preamble={\usetikzlibrary{3dtools}}]
\begin{tikzpicture}[declare function={a=1.5;},
   dot/.style={circle,inner sep=1pt,fill},
   3d/install view={phi=100,psi=0,theta=70}]	
  \path (a,0,0) coordinate (A)
   (0,a,0) coordinate (B)
   (0,0,a) coordinate (C);
  \path[3d/intersection of three spheres=
  	{rA=2,rB=2.5,rC=3}];
  \path (i1) 
    coordinate[dot,label=above:{$%
  	\pgfmathparse{TD("(i1)")}%
   (\pgfmathprintvector\pgfmathresult)^T$}];
  \path (i2) 
    coordinate[dot,label=above:{$%
  	\pgfmathparse{TD("(i2)")}%
   (\pgfmathprintvector\pgfmathresult)^T$}];
  \path foreach \X in {A,B,C}
  {(\X) coordinate[dot,label=above:{$\X$}]};
\end{tikzpicture}
\end{codeexample}

\begin{key}{/tikz/3d/insphere center=\meta{options} (initially empty)}
		Computes the center of an insphere center of a tetrahedron defined by
		four corners, |A|, |B|, |C|, and |D|. The initial values are
		|A=(A)|, |B=(B)|, |C=(C)|, and |D=(D)|. 
\end{key}

\begin{codeexample}[width=5.4cm,preamble={\usetikzlibrary{3dtools}}]
\begin{tikzpicture}[line join=round,line cap=round,
	3d/install view={phi=100,psi=0,theta=70}]	
   \path (0,0,0) coordinate (D)
	(2,0,0) coordinate (A)
	(0,3,0) coordinate (B)
	(0,0,4) coordinate (C);
   \path[3d/insphere center] coordinate (I);	
   \tikzset{3d/plane through=(A) and (B) and (C) named pABC}
   \path[3d/project={(I) on pABC}] coordinate (I');
   \pgfmathsetmacro{\myr}{tddistance("(I)","(I')")}
   \draw (A) -- (B) -- (C) -- cycle --(D) --(B) (C) -- (D);
   \shade[3d/screen coords,ball color=green, opacity=0.8] 
		(I) circle [radius=\myr];
\end{tikzpicture}
\end{codeexample}



\begin{key}{/tikz/pics/3d circle through 3 points=\meta{options} (initially empty)}
        Draws a circle through 3 points in 3 dimensions. If the three
		coordinates are close to linearly dependent, the circle will not be
		drawn. \emph{This pic will most likely be removed from the documentation
		and no longer be supported/developed.}
\end{key}
\begin{key}{/tikz/3d circle through 3 points/A (initially {(1,0,0)})}
        First coordinate. Can be either symbolic or explicit. Symbolic
		coordinates need to be defined via 
		|\path (x,y,z) coordinate (name);|.
\end{key}
\begin{key}{/tikz/3d circle through 3 points/B (initially {(0,1,0)})}
        Second coordinate, like above.
\end{key}
\begin{key}{/tikz/3d circle through 3 points/C (initially {(0,0,1)})}
        Third coordinate, like above.
\end{key}
\begin{key}{/tikz/3d circle through 3 points/center name (initially {M})}
        Name of the center coordinate that will be derived.
\end{key}
\begin{key}{/tikz/3d circle through 3 points/auxiliary coordinate prefix (initially {tmp})}
		In \tikzname\ the coordinates are global. The code for the circle is
		more comprehensible if named coordinates are introduced. Their names
		will begin with this prefix. Changing the prefix will allow users to
		avoid overwriting existing coordinates.
\end{key}

\begin{codeexample}[width=4.5cm,preamble={\usetikzlibrary{3dtools}}]
\begin{tikzpicture}[3d/install view={%
    phi=30,psi=0,theta=70}]
 \foreach \X in {A,B,C}
 {\pgfmathsetmacro{\myx}{3*(rnd-1/2)}
 \pgfmathsetmacro{\myy}{3*(rnd-1/2)}
 \pgfmathsetmacro{\myz}{3*(rnd-1/2)}
 \path (\myx,\myy,\myz) coordinate (\X);}
 \path pic{3d circle through 3 points={%
 A={(A)},B={(B)},C={(C)},center name=MM}};
  \foreach \X in {A,B,C,MM}
  {\fill (\X) circle[radius=1.5pt] 
  node[above]{$\X$};}
\end{tikzpicture}
\end{codeexample}

\begin{key}{/tikz/pics/3d/circle on sphere=\meta{options} (initially empty)}
		Draws a circle (or an arc) on a sphere. It distinguishes between
		foreground and background. Notice that this pic defines its own
		dreibein, so it will define coordinates that draw their names from
		|/tikz/3d/aux keys/ex|, |/tikz/3d/aux keys/ey| and
		|/tikz/3d/aux keys/ez|. It will also define a coordinate |(nscreen)|
		which contains the normal to the screen.  The foreground path will use
		the style |/tikz/3d/visible|, and the hidden path will be drawn with the
		style |/tikz/3d/hidden|. The commands of the
		\href{https://github.com/matthias-wolff/tikz-3dplot-circleofsphere}{%
		\texttt{tikz-3dplot-circleofsphere}} package are available through the
		library |3dtools.circleofsphere| (section~\ref{sec:CircleOfSphere}).
\end{key}
\begin{key}{/tikz/3d/circle on sphere/C (initially {(0,0,0)})}
        Coordinate of the center of the sphere.
\end{key}
\begin{key}{/tikz/3d/circle on sphere/P (initially {(0,0,0)})}
        Coordinate of the center of the circle.
\end{key}
\begin{key}{/tikz/3d/circle on sphere/n (initially {(0,0,1)})}
        Normal. This is only needed if the center of the circle coincides with
		the center of the sphere.
\end{key}
\begin{key}{/tikz/3d/circle on sphere/auxiliary coordinate prefix (initially {tmp})}
        This pic installs a symbolic coordinate for the center of the circle.
		The auxiliary prefix can be used to avoid overwriting existing
		coordinates. The initial choice of the prefix will lead to a coordinate
		with name |(tmpP)|.
\end{key}
\begin{key}{/tikz/3d/circle on sphere/fore layer=\meta{layer} (initially main)}
 		Layer of the visible faces.
\end{key}

\begin{key}{/tikz/3d/circle on sphere/back layer=\meta{layer} (initially main)}
 		Layer of the hidden faces.
\end{key}
\begin{key}{/tikz/3d/circle on sphere/latitude=\meta{angle} (initially empty)}
		The latitude circle at \meta{angle}: sets |n| to $(0,0,1)$ and |P| to
		|C| shifted by $R\sin(\text{\meta{angle}})$ along $z$.
\end{key}
\begin{key}{/tikz/3d/circle on sphere/longitude=\meta{angle} (initially empty)}
		The meridian at the azimuth \meta{angle}: the great circle through the
		poles and the direction $(\cos\text{\meta{angle}},\sin\text{\meta{angle}},0)$.
\end{key}
\begin{key}{/tikz/3d/circle on sphere/fill (initially empty)}
		Style of the disk of the circle, which is painted first, e.g.\
		|fill/.style={fill=red,fill opacity=0.2}|.
\end{key}
\begin{key}{/tikz/3d/circle on sphere/start angle=\meta{angle} (initially 0)}
		Start of the arc. Angles are measured in the plane of the circle from
		the projection of |u|, positively around |n|. The arc gets split where
		it passes from the front of the sphere to the back.
\end{key}
\begin{key}{/tikz/3d/circle on sphere/end angle=\meta{angle} (initially 360)}
		End of the arc.
\end{key}
\begin{key}{/tikz/3d/circle on sphere/u (initially {(1,0,0)})}
		Direction of the angle 0 (it gets projected into the plane of the circle).
\end{key}
\begin{key}{/tikz/3d/circle on sphere/start color=\meta{color} (initially empty)}
		If set, the arc gets a color gradient from \meta{color} at its start to
		|end color| at its end. It is drawn as sub-arcs of at most
		|segment angle| degrees, each with its own color and, in a scene, its own
		depth. Only the last sub-arc keeps the arrow tips, so put an end tip
		(e.g.\ |->|) in the styles of the parts.
\end{key}
\begin{key}{/tikz/3d/circle on sphere/end color=\meta{color} (initially empty)}
		Color at the end of the arc (initially the start color).
\end{key}
\begin{key}{/tikz/3d/circle on sphere/segment angle=\meta{angle} (initially 360)}
		In a scene, the arc is painted as cells of at most this angle (on a
		grid that starts at the point closest to the viewer), each with the
		depth of its centroid. Sub-arcs of 1pt or more get round caps, which
		close their joints.
\end{key}
\begin{key}{/tikz/3d/circle on sphere/gradient angle=\meta{angle} (initially 2)}
		Largest angle of a sub-arc of a color gradient.
\end{key}

\begin{codeexample}[width=4.5cm,preamble={\usetikzlibrary{3dtools}}]
\begin{tikzpicture}[3d/install view={phi=110,theta=70},
	declare function={R=2;}]
 \shade[ball color=white,3d/screen coords] circle[radius=R];
 \path pic{3d/circle on sphere={R=R,P={(0.5,-1.2,1)}}};
\end{tikzpicture}
\end{codeexample}

\begin{codeexample}[width=4.5cm,preamble={\usetikzlibrary{3dtools}}]
\begin{tikzpicture}[3d/install view={phi=110,theta=70},
	declare function={R=2;}]
 \shade[ball color=white,3d/screen coords,opacity=0.5]
   circle[radius=R];
 \path pic{3d/circle on sphere={R=R,latitude=30,
     fill/.style={fill=red,fill opacity=0.2}}}
   pic{3d/circle on sphere={R=R,longitude=40}}
   pic{3d/circle on sphere={R=R,n={(1,1,0)},end angle=270,
     start color=blue,end color=red,
     /tikz/3d/visible/.style={thick,->},
     /tikz/3d/hidden/.style={thick,->,opacity=0.4}}};
\end{tikzpicture}
\end{codeexample}

\begin{key}{/tikz/pics/3d/sphere=\meta{options} (initially empty)}
		A sphere of radius |R| around the pic position, painted with the style
		|shell| and outlined with |outline|. In a scene the shell is painted
		twice, behind and in front of everything inside the sphere, and the
		visible parts of |3d/circle on sphere| go on top of it; for a
		translucent sphere, the opacity of |shell| applies to each layer.
\end{key}
\begin{key}{/tikz/3d/sphere/R (initially 1)}
		Radius.
\end{key}
\begin{key}{/tikz/3d/sphere/shell (initially {ball color=white})}
		Style of the shell.
\end{key}
\begin{key}{/tikz/3d/sphere/outline (initially draw)}
		Style of the outline.
\end{key}

\begin{key}{/tikz/pics/3d/point on sphere=\meta{options} (initially empty)}
		Puts a node at the point |P| on the sphere around |C|, with the style
		|front| and the text |front text| if the point faces the viewer, and
		with |back| and |back text| otherwise.
\end{key}
\begin{key}{/tikz/3d/point on sphere/C (initially {(0,0,0)})}
		Center of the sphere.
\end{key}
\begin{key}{/tikz/3d/point on sphere/P (initially {(0,0,1)})}
		The point.
\end{key}
\begin{key}{/tikz/3d/point on sphere/front text (initially \textbackslash textbullet)}
		Text of the node on the front.
\end{key}
\begin{key}{/tikz/3d/point on sphere/back text (initially \textbackslash(\textbackslash circ\textbackslash))}
		Text of the node on the back.
\end{key}

\begin{key}{/tikz/pics/3d/circle=\meta{options} (initially empty)}
		A circle (or an arc) around the position of the pic, split into the
		half that is closer to the viewer, drawn with |front|, and the other
		half, drawn with |back|. In a scene (section~\ref{sec:Scenes}) the
		halves are separate pieces, so the circle can wrap around other
		objects. The keys |n|, |u|, |start angle|, |end angle|, |fill|,
		|start color|, |end color|, |segment angle| and |gradient angle| work
		as for |3d/circle on sphere|.
\end{key}
\begin{key}{/tikz/3d/circle/r (initially 1)}
		Radius.
\end{key}
\begin{key}{/tikz/3d/circle/front (initially empty)}
		Style of the half closer to the viewer.
\end{key}
\begin{key}{/tikz/3d/circle/back (initially empty)}
		Style of the other half.
\end{key}

\begin{key}{/tikz/pics/3d incircle=\meta{options} (initially empty)}
        Inscribes a circle in a triangle in 3 dimensions.
\end{key}
\begin{key}{/tikz/3d incircle/A (initially {(1,0,0)})}
        First coordinate. Can be either symbolic or explicit. 
\end{key}
\begin{key}{/tikz/3d incircle/B (initially {(0,1,0)})}
        Second coordinate, like above.
\end{key}
\begin{key}{/tikz/3d incircle/C (initially {(0,0,1)})}
        Third coordinate, like above.
\end{key}
\begin{key}{/tikz/3d incircle/center name (initially {I})}
        Name of the center coordinate that will be derived.
\end{key}
\begin{key}{/tikz/3d incircle/projection prefix (initially {tmpp})}
		In \tikzname\ the coordinates are global. The code for the circle is
		more comprehensible if named coordinates are introduced. Their names
		will begin with this prefix. Changing the prefix allows users to
		avoid overwriting existing coordinates.
\end{key}

\begin{codeexample}[width=2.5cm,preamble={\usetikzlibrary{3dtools}}]
\begin{tikzpicture}[3d/install view={phi=110,psi=0,theta=70}]
 \draw
   (8,5,5) coordinate[label=above:{$A$}] (A) --
   (1,2,0) coordinate[label=below:{$B$}] (B) --
   (5,-5,0) coordinate[label=below:{$C$}] (C) -- cycle;
 \path pic[red,dashed]{3d incircle={% 
 A={(A)},B={(B)},C={(C)},center name=I}};  
 \draw (I) -- (tmppa) (I) -- (tmppb) (I)  -- (tmppc);
\end{tikzpicture}
\end{codeexample}

\begin{key}{/tikz/3d/polyhedron/draw face with corners=\meta{list of corners} (initially empty)}
		Draws a face of a polyhedron through the specified corners. 
		It distinguishes between ``visible'' and ``hidden'' faces. The user can
		specify a point inside the corner. Then the outwards pointing normal of
		the face can either have a negative projection on the screen, in which
		case the face is hidden, or else it is visible.
\end{key}


\begin{key}{/tikz/3d/polyhedron/draw open face with corners=\meta{list of corners} (initially empty)}
		Same as the previous key except it does not close the surface boundary
		path.
\end{key}

\begin{key}{/tikz/3d/polyhedron/O=\meta{coordinate} (initially {(0,0,0)})}
 		Coordinate inside the polyhedron.
\end{key}

\begin{key}{/tikz/3d/polyhedron/L=\meta{vector} (initially the value of |3d/light|)}
 		Direction towards the light, used for the faces of polyhedra only.
		By default it follows |/tikz/3d/light|.
\end{key}

\begin{key}{/tikz/3d/polyhedron/shading function=\meta{name} (initially tikztdpolyhedronshade)}
 		Function of the cosine between the outward normal of the face and
		the light direction. It returns the percentage of the face color
		that is mixed with black. The default,
		|tikztdpolyhedronshade(\x)=100*tikztdlight(\x)|, uses the lighting
		model below. The look of earlier versions is obtained with
		|declare function={tikztdpolyhedronshade(\x)=70+30*\x;}|.
\end{key}

\begin{key}{/tikz/3d/light=\meta{vector} (initially {(1,1,1)})}
		Direction pointing towards a single directional light, in the
		coordinates in which the drawing is specified (the frame of
		|3d/install view|). The light is thus fixed in the scene: changing the
		view changes what is lit as seen from the camera. It is used by
		polyhedra, |cylinder|, |ycylinder| and |3d/shaded cone|. A surface
		element with unit outward normal $n$ and the unit light direction $L$
		gets the intensity
		$I=\min\bigl(1,a+d\max(0,n\cdot L)\bigr)$ (Lambert), and the color
		\meta{color}|!|$100I$|!black|. The pgf function |tikztdlight| computes
		$I$ from $n\cdot L$.
\end{key}
\begin{key}{/tikz/3d/light/ambient=\meta{number} (initially 0.5)}
		The ambient part $a$.
\end{key}
\begin{key}{/tikz/3d/light/diffuse=\meta{number} (initially 0.5)}
		The diffuse part $d$.
\end{key}


\begin{key}{/tikz/3d/polyhedron/fore layer=\meta{layer} (initially main)}
 		Layer of the visible faces.
\end{key}

\begin{key}{/tikz/3d/polyhedron/back layer=\meta{layer} (initially main)}
 		Layer of the hidden faces.
\end{key}

\begin{key}{/tikz/3d/polyhedron/color=\meta{color} (initially yellow)}
 		Color of the polyhedron.
\end{key}

\begin{key}{/tikz/3d/polyhedron/fore=\meta{style} (initially {draw,solid})}
 		Style for visible faces. One can turn of the fill by adding |fill=none|.
\end{key}


\begin{key}{/tikz/3d/polyhedron/back=\meta{style} (initially {draw,dashed,fill opacity=0})}
 		Style for visible faces.
\end{key}

\begin{key}{/tikz/3d/polyhedron/complete dashes=\meta{style} (initially {on 2pt off 2pt})}
 	    Switches on dashes for the edges which start and end with an 
		on phase.
\end{key}


\begin{key}{/tikz/cheating dash=\meta{style} (initially {on 2pt off 2pt})}
 	    Switches on dashes with an on phase.
		Taken from \href{https://tex.stackexchange.com/a/133357}{Mark Wibrow's answer}.
\end{key}

\begin{codeexample}[width=4.5cm,preamble={\usetikzlibrary{3dtools}}]
\pgfdeclarelayer{background} 
\pgfdeclarelayer{foreground}
\pgfsetlayers{background,main,foreground}
\begin{tikzpicture}
\begin{scope}[3d/install view={phi=110,psi=10,%
	theta=70},scale=2]
  \path (0,0,1) coordinate (v1) 
	  ({sqrt(8/9)},0,-1/3) coordinate (v2)
	  ({-sqrt(2/9)},{sqrt(2/3)},-1/3) coordinate (v3)
	  ({-sqrt(2/9)},{-sqrt(2/3)},-1/3) coordinate (v4);
  \tikzset{3d/polyhedron/.cd,
	fore layer=foreground,back layer=background,
	fore/.append style={opacity=0.4,thick},
	back/.append style={3d/polyhedron/complete dashes,
		opacity=0.4},
 	draw face with corners={{(v1)},{(v2)},{(v3)}},
 	draw face with corners={{(v1)},{(v2)},{(v4)}},
 	draw face with corners={{(v1)},{(v3)},{(v4)}},
 	draw face with corners={{(v2)},{(v3)},{(v4)}}}	  
\end{scope}
\end{tikzpicture}
\end{codeexample}


\begin{key}{/tikz/3d/define vertices=\meta{list} (initially empty)}
   Allows the user to define a list of vertices from a list. These vertices have
   names that consist of their indices in the list. One can work with the 
   |name prefix| key to give them more meaningful names. Such lists can be generated
   with \texttt{Mathematica}, e.g.\ with
   \texttt{N[PolyhedronData["Dodecahedron", "VertexCoordinates"]]}.
\end{key}


\begin{key}{/tikz/3d/polyhedron/create faces from vertex list=\meta{list} (initially empty)}
   Draws a face from a list of vertices. Such lists can be generated with  with
   \texttt{Mathematica}, e.g.\ with \texttt{PolyhedronData["Dodecahedron",
   "FaceIndices"]}. This key is very similar to  |draw face with corners| except
   that it also adds the |name prefix|, and that it loops over the list.
   One should avoid spaces in the list.
\end{key}


\begin{codeexample}[width=3.5cm,preamble={\usetikzlibrary{3dtools}}]
\pgfdeclarelayer{background}\pgfdeclarelayer{foreground}
\pgfsetlayers{background,main,foreground}
\begin{tikzpicture}
 \edef\lstV{{-1.376,0.,0.2628},{1.376,0.,-0.2628},{-0.4253,-1.309,0.2628},
 {-0.4253,1.309,0.2628},{1.113,-0.8090,0.2628},{1.113,0.8090,0.2628},
 {-0.2628,-0.8090,1.113},{-0.2628,0.8090,1.113},{-0.6881,-0.5,-1.113},
 {-0.6881,0.5,-1.113},{0.6881,-0.5,1.113},{0.6881,0.5,1.113},
 {0.8506,0.,-1.113},{-1.113,-0.8090,-0.2628},{-1.113,0.8090,-0.2628},
 {-0.8506,0.,1.113},{0.2628,-0.8090,-1.113},{0.2628,0.8090,-1.113},
 {0.4253,-1.309,-0.2628},{0.4253,1.309,-0.2628}}
 \edef\lstFaces{{15,10,9,14,1},{2,6,12,11,5},{5,11,7,3,19},{11,12,8,16,7},
 {12,6,20,4,8},{6,2,13,18,20},{2,5,19,17,13},{4,20,18,10,15},{18,13,17,9,10},
 {17,19,3,14,9},{3,7,16,1,14},{16,8,4,15,1}}
 \tikzset{3d/polyhedron/.cd,fore layer=foreground,back layer=background, 
	fore/.append style={fill opacity=0.7},
	back/.append style={3d/polyhedron/complete dashes,fill opacity=0.7}}
\begin{scope}[3d/install view={phi=120,psi=20,%
	theta=70}]
 \tikzset{name prefix=Va,%<- used for all vertices
 	3d/define vertices/.expanded={\lstV},
	3d/polyhedron/create faces from vertex list/.expanded={\lstFaces}}
\end{scope}
\begin{scope}[yshift=-3.5cm,
	3d/install view={phi=210,psi=30,theta=60}]
 \tikzset{name prefix=Vb,%<- used for all vertices
 	3d/define vertices/.expanded={\lstV},
 	3d/polyhedron/create faces from vertex list/.expanded={\lstFaces}}
\end{scope}
\begin{scope}[yshift=-7cm,
	3d/install view={phi=30,psi=-10,theta=75}]
 \tikzset{name prefix=Vc,%<- used for all vertices
 	3d/define vertices/.expanded={\lstV},
 	3d/polyhedron/create faces from vertex list/.expanded={\lstFaces}}
\end{scope}
\end{tikzpicture}
\end{codeexample}

The library has also some experimental pics for cones, truncated cones and
cylinders. They share a couple of common keys. 

\begin{key}{/tikz/3d/r (initially 1)}
 Key for radii, e.g.\ of cylinders or cones.
\end{key}
\begin{key}{/tikz/3d/h (initially 1)}
 Key for heights,  e.g.\ of cylinders or cones.
\end{key}
\begin{key}{/tikz/3d/R (initially 2)}
 Key for larger radii, e.g.\ of a frustum.
\end{key}
\begin{key}{/tikz/pics/3d/visible (initially {draw,solid})}
 Style for visible lines.
\end{key}
\begin{key}{/tikz/pics/3d/hidden (initially {draw,very thin,cheating dash})}
 Style for hidden lines.
\end{key}

\begin{key}{/tikz/pics/3d/cone (initially empty)}
 Draws a cone whose base, of radius |r|, is centered at the pic position and
 whose apex sits at distance |h| along |3d/cone/axis|. The pic defines the
 coordinates |-base|, |-apex| and, for compatibility, |-tip| (same as |-apex|).
 Hidden parts of the base rim are drawn with |3d/hidden|.
\end{key}
\begin{key}{/tikz/3d/cone/axis (initially {(0,0,1)})}
 Direction of the cone axis (it gets normalized). Also used by |3d/frustum|.
\end{key}
\begin{key}{/tikz/3d/cone/apex=\meta{3d vector}}
 Apex relative to the base center. Sets |3d/cone/axis| and |3d/h|.
\end{key}
\begin{key}{/tikz/3d/cone/outer (initially empty)}
 Style of the visible part of the mantle, e.g.\ |fill=red!30|.
\end{key}
\begin{key}{/tikz/3d/cone/inner (initially empty)}
 Style of the base region when the base side faces the viewer (the base disk,
 or the interior of an open cone).
\end{key}
\begin{key}{/tikz/3d/cone/outer shade (initially 3d/cone/lit side=outer)}
 Shading of the mantle used by |3d/shaded cone|. It is centered at the apex.
 By default the mantle is lit by |3d/light|.
\end{key}
\begin{key}{/tikz/3d/cone/inner shade (initially 3d/cone/lit side=inner)}
 Shading of the base region (the interior of the open cone) used by
 |3d/shaded cone|, lit by |3d/light|.
\end{key}
\begin{key}{/tikz/3d/cone/color=\meta{color} (initially white)}
 Color of the lit cone.
\end{key}
\begin{key}{/tikz/3d/cone/lit side=\meta{outer or inner}}
 Sets |3d/cone/lit/max| and |3d/cone/lit/min| to the intensities of the
 brightest and darkest visible generator of that side and calls
 |3d/cone/light shading| with the screen angle, seen from the apex, of the
 brightest generator.
\end{key}
\begin{key}{/tikz/3d/cone/light shading=\meta{angle}}
 The shading hook. The default declares a functional shading in the angle
 around the apex that reproduces the exact Lambert intensity of every visible
 generator (brightness is constant along a generator). Replace it to use
 a different shading, e.g.\ one built from \meta{angle} and the two keys above.
\end{key}
\begin{key}{/tikz/3d/cone/screen angle}
 Set by the cone pics before the shading styles are used: the screen angle of
 the direction from the base center to the apex. For instance,
 |outer shade/.style={3d/cone shading={alpha/.evaluated={\pgfkeysvalueof{/tikz/3d/cone/screen angle}-90}}}|
 rotates the shading with the cone.
\end{key}

\begin{key}{/tikz/pics/3d/shaded cone (initially empty)}
 Draws a shaded cone with the basis in the current xy plane.
\end{key}

\begin{key}{/tikz/3d/cone shading=\meta{options}}
 Defines and uses a functional shading that can be used for cones. The
 parameters are the coordinates of the tip |x| and |y|, which have to satisfy
 $-25<x,y<25$, and |alpha|, which indicates the rotation of the shading.
\end{key}


\begin{codeexample}[width=4.8cm,preamble={\usetikzlibrary{calc,3dtools}}]
\begin{tikzpicture}[3d/install view=%
	{phi=110,psi=0,theta=60}]
 \pic{3d/shaded cone={r=2,h=3}};
\end{tikzpicture}
\end{codeexample}



\begin{key}{/tikz/pics/3d/frustum (initially empty)}
 Draws a frustum with the basis in the current xy plane, open at both ends.
 The generatrices of the silhouette divide the mantle in two patches: one
 shows its outer side to the viewer, painted with |3d/frustum/outer|, the
 other its inner side, painted with |3d/frustum/inner|. The rims and the
 generatrices are drawn with |3d/visible| and |3d/hidden|. Notice that if
 both radii are equal this is just a cylinder.
\end{key}
\begin{key}{/tikz/3d/frustum/outer (initially empty)}
 Style of the patch of the mantle that shows its outer side.
\end{key}
\begin{key}{/tikz/3d/frustum/inner (initially empty)}
 Style of the patch of the mantle that shows its inner side, e.g.\ the
 coating of a conical mirror.
\end{key}


\begin{codeexample}[width=4.8cm,preamble={\usetikzlibrary{calc,3dtools}}]
\begin{tikzpicture}[3d/install view=%
	{phi=110,psi=0,theta=60}]
 \pic{3d/frustum};
\end{tikzpicture}
\end{codeexample}



\begin{key}{/tikz/pics/ycylinder (initially empty)}
 A cylinder in the $y$-direction. This pic requires the |calc| library. 
 As of now it does only work for |psi=0|. 
 This key is obsolete because of the |cylinder| pic which allows the user to dial an arbitrary axis.
\end{key}
\begin{key}{/tikz/pics/cylinder (initially empty)}
	A cylinder with an arbitrary axis. 
	This pic requires the |calc| library. 
	Seen along its axis it reduces to the cap facing the viewer.
\end{key}
   
\begin{key}{/tikz/pics/3d/cylinder/axis (initially {(0,0,1)})}
	Axis of cylinder.
\end{key}
   
\begin{key}{/tikz/pics/3d/cylinder/mantle (initially draw)}
 Style for cylinder mantle. If it contains no fill or shading, the mantle
 shows the lighting of |3d/cylinder/lighting|. A user shading (e.g.\
 |left color|, |right color|, |middle color|) is oriented across the axis.
\end{key}
\begin{key}{/tikz/3d/cylinder/color=\meta{color} (initially white)}
 Color of the lit mantle and caps.
\end{key}
\begin{key}{/tikz/3d/cylinder/lighting (initially a path picture with the lit shading)}
 Painted below the mantle and the cap styles. In parallel projection the
 Lambert intensity of the mantle only depends on the position across the
 silhouette, so the default axial shading is exact. A cap is flat: the
 visible one, with outward normal $n=\pm$axis, gets the constant intensity
 $\min\bigl(1,a+d\max(0,n\cdot L)\bigr)$. It is not painted below a
 mantle or cap whose style fills or shades. Set it to |{}| to switch the
 lighting of mantle and caps off.
\end{key}
\begin{key}{/tikz/3d/cylinder/hole=\meta{radius} (initially 0)}
	Radius of a hole in the caps.
\end{key}
\begin{key}{/tikz/3d/cylinder/open=\meta{boolean} (initially false)}
	A tube without caps: only the rim of the cap facing the viewer is drawn.
\end{key}
\begin{key}{/tikz/pics/3d/cylinder/top (initially draw)}
 Style for the visible cap (top or bottom, also used by |ycylinder|). If it
 contains no fill, the cap shows the lighting of |3d/cylinder/lighting|.
\end{key}



\begin{codeexample}[width=3cm,preamble={\usetikzlibrary{3dtools}}]
\begin{tikzpicture}[3d/install view={phi=30,psi=0,theta=80}]
 \pic{cylinder={r=0.4,h=3,
 cylinder/axis={(1,0,1)},
 cylinder/top/.style={fill=blue}}};
\end{tikzpicture}
\end{codeexample}

Seen along its axis, a cylinder shows only the cap facing the viewer, with its
hole, or only the rim if it is open.
\begin{codeexample}[width=5cm,preamble={\usetikzlibrary{3dtools}}]
\begin{tikzpicture}[3d/install view={phi=0,psi=0,theta=0}]
 \fill[red] (-1.2,-0.15) rectangle (4.2,0.15);
 \pic{cylinder={r=1,h=2,cylinder/hole=0.5,
  cylinder/top/.style={fill=blue!50,draw}}};
 \pic at (3,0,0) {cylinder={r=1,h=2,cylinder/open=true}};
\end{tikzpicture}
\end{codeexample}

\begin{codeexample}[width=5cm,preamble={\usetikzlibrary{calc,3dtools}}]
\begin{tikzpicture}[3d/install view={phi=110,psi=0,theta=70},
  3d/light={(-1,0.3,0.5)},3d/cylinder/color=teal!60,
  3d/cone/color=orange]
 \pic{cylinder={r=0.6,h=2}};
 \begin{scope}[xshift=2.5cm]
  \pic{3d/shaded cone={r=1,h=2}};
 \end{scope}
\end{tikzpicture}
\end{codeexample}

\begin{key}{/tikz/pics/3d/dim line (initially empty)}
	Draws a line that allows users to annotate axes and/or indicate dimensions.
	It expects the user to specify the start and end points, |A| and |B|, as well as the normal |n| of the plane the line is to be drawn. 
	In addition, the |d| key, which can take negative values, determines the distance of the line from |A--B|.
	The label is given by the |l| key, and  
	the |dim label| style controls its appearance.
	This |pic| requires the |calc| library. 
\end{key}

\begin{codeexample}[width=6cm,preamble={\usetikzlibrary{calc,3dtools}}]
\begin{tikzpicture}[3d/install view={phi=110,theta=60},
	declare function={a=4;rx=2.5;ry=3;rz=2.5;}]
 \draw[gray,help lines] (0,0,0) coordinate (O) 
 	-- (a,0,0)  
	(O) -- (0,a,0) (O) -- (0,0,a);
 \path (rx,ry,rz) coordinate[circle,inner sep=1.5pt,
			fill] (r);
 \path pic{3d/dim line={A={(0,ry,0)},B={(rx,ry,0)},l={$x$}}}
			pic{3d/dim line={A={(rx,0,0)},B={(rx,ry,0)},l={$y$},d=-2mm}}
			pic{3d/dim line={A={(rx,ry,0)},B={(r)},n={(1,0,0)},l={$z$},
			d=-2mm,dim label/.append style={rotate=90}}};
\end{tikzpicture}
\end{codeexample}

% \begin{key}{/tikz/3d/cone/r (initially 1)}
%  Radius of the base of the cone.
% \end{key}
% \begin{key}{/tikz/3d/cone/h (initially 1)}
%  Height of the cone.
% \end{key}
% \begin{key}{/tikz/pics/3d/h (initially 1)}
%  Key for heights,  e.g.\ of cylinders.
% \end{key}
% \begin{key}{/tikz/pics/3d/mantle (initially draw)}
%  Style for cylinder mantle. If no fill option is specified, it will be shaded.
% \end{key}
% \begin{key}{/tikz/pics/3d/top (initially draw)}
%  Style for cylinder top. 
% \end{key}

To do:
\begin{itemize}
 \item transform to plane given by three non-degenerate coordinates
 \item transform to plane given by normal and one point
 \item maybe layering/visibility
\end{itemize}

\subsection{Path management}

This is poor marmots version of path management. Almost certainly
\href{https://ctan.org/pkg/spath3?lang=en}{\texttt{spath3}} is way superior. The
idea is that one can store paths, and redraw them, possibly in reversed order.
This recording tool uses the |show path construction| decoration from the
|decorations.pathreplacing| library. In order to make use of these tools, you
need to load the library.

\begin{key}{/tikz/store path=\meta{name} (initially empty)}
 		Records a path under the name \meta{name}. If you want to draw and
		record a path at the same time, use |postaction={store path=name}|.
\end{key}


\begin{key}{/tikz/stored path/coordinate prefix=\meta{prefix} (initially stored-)}
 		Prefix for the stored coordinates. If you store several paths and want
		to possibly combine them, you may want to give each path a unique
		prefix.
\end{key}

\begin{key}{/tikz/stored path/reset coordinate index}
 		Resets the counter for the stored coordinates. Mainly for internal use.
\end{key}


\begin{key}{/tikz/stored path/restore path=\meta{name} (initially empty)}
 		Restores a saved path of name \meta{name} using the |insert path| key.
\end{key}

\begin{key}{/tikz/stored path/restore reversed path=\meta{name} (initially empty)}
 		Restores a saved path of name \meta{name} using the |insert path| key
		in \emph{reversed} order.
\end{key}


\begin{key}{/tikz/stored path/append path=\meta{name} (initially empty)}
		Restores a saved path of name \meta{name} using the |insert path| key with a
		|--| at the beginning. This may be necessary to patch paths together.
\end{key}

\begin{key}{/tikz/stored path/append reversed path=\meta{name} (initially empty)}
		Restores a saved path of name \meta{name} using the |insert path| key in
		\emph{reversed} order with a |--| at the beginning. This may be necessary to
		patch paths together.
\end{key}

\begin{key}{/tikz/stored path/first coordinate of=\meta{name} (initially empty)}
 		Inserts the first coordinate of the path named \meta{name} using the |insert path| key.
\end{key}

\begin{key}{/tikz/stored path/last coordinate of=\meta{name} (initially empty)}
 		Inserts the last coordinate of the path named \meta{name} using the |insert path| key.
\end{key}

\begin{key}{/tikz/stored path/coordinate with index=\mvar{n} of \meta{name} (initially empty)}
 		Inserts the \mvar{n}\textsuperscript{th} coordinate of the path named \meta{name} using the |insert path| key.
		Notice that at present there is no sanity check. If \mvar{n} is
		larger than the number of stored coordinates, there will be an error.
\end{key}

\begin{math-function}{tikztdindexoflastcoordinate}
   Yields index of the last coordinate of a stored path. This can be used to
   test if a path is empty.
\end{math-function}


\subsection{List management}

3d ordering has a lot to do with sorting lists. Here are some poor marmot's
tools to deal with such lists and some primitive versions of associative arrays.
First of all, there are a couple of \emph{key handlers} that allow one to add
and remove items. Notice that, in order to apply key handlers to a list, the
list has to be stored in a \meta{key}.

\begin{handler}{{.add item}}
		Key handler that adds an item to a list \meta{key}.
\end{handler}

\begin{handler}{{.remove item}}
		Key handler that removes an item from a list \meta{key}.
\end{handler}

\begin{handler}{{.count}|=|\meta{macro}}
		Key handler that stores the number of elements of a given list \meta{key} in
		a macro, which is the mandatory argument.
\end{handler}

\begin{handler}{{.print list}}
		Key handler that prints a list. It threads 
		|/list management/print item|
		code over the list, and separates the items by 
		|/list management/list separators|.
\end{handler}

\begin{key}{/list management/print item}
 		Code used by the |.print list| key handler.
\end{key}

\begin{key}{/list management/list separator (initially {,})}
 		Code used by the |.print list| key handler to separate items.
\end{key}


\begin{codeexample}[width=6.8cm,preamble={\usetikzlibrary{3dtools}}]
\begin{minipage}{6.4cm}
\pgfkeys{/my lists/.cd,
tikzlings/.initial={anteater,%
bear,bee,cat,coati,%
hippo,koala,marmot,mouse}}
The first\ 
\pgfkeys{/my lists/tikzlings/.count=\myn}\myn\ 
Ti\emph{k}Zlings are\ 
\{\pgfkeys{/my lists/tikzlings/.print list}\}.
\par\medskip
\pgfkeys{/my lists/tikzlings/.add item=wolpertinger}
After the wolpertinger joined, they were \pgfkeys{/my
lists/tikzlings/.count=\myn}{$\myn$}, and the list became\ 
\{\pgfkeys{/my lists/tikzlings/.print list}\}.
\par\medskip
\pgfkeys{/my lists/tikzlings/.remove item=mouse}
Sadly, the mouse could not stand the wolpertinger,\ 
so the list Ti\emph{k}Zlings got reduced to\ 
\pgfkeys{/my lists/tikzlings/.count=\myn}{$\myn$},\ 
consisting of\
\{\pgfkeys{/my lists/tikzlings/.print list}\}.
\end{minipage}
\end{codeexample}

\begin{handler}{{.is array}|=|\meta{list}}
		Key handler makes the current key \meta{key} an array consisting of
		\meta{list}. You need to use a new key for that, i.e.\ you cannot
		overwrite existing keys. The benefit is that you can access the items in
		the list simply via $\meta{key}/n$, where $n$ is an integer-valued
		index, and that \texttt{\meta{key}/dim} contains its dimension.
\end{handler}

\begin{handler}{{.sort numeric list}|=|{\meta{macro}}{\meta{macro}}}
		Key handler that sorts a list \meta{key}. The resulting sorted list is
		stored in the first argument, which has to be a macro, and the
		permutation is stored in the second argument, which needs to be a macro
		as well. This allows one to know the position of a given entry in a
		sorted list.
\end{handler}

\begin{codeexample}[width=7cm,preamble={\usetikzlibrary{3dtools}}]
\begin{minipage}{6.4cm}
\pgfkeys{/my lists/.cd,
my initial array/.is array={%
	12,19,1,3,17,4,23,5,9,7}, 
my values/.initial=\pgfkeysvalueof{%
	/my lists/my initial array/content},%
my values/.sort numeric list=%
{\temp}{\templ},% 
my sorted array/.is array/.expanded={\temp}, 
my index machinery/.is array/.expanded={\templ}}% 
Often one is given a list of entries where\ 
one wants to sort on the basis of some\ 
function that can be fed with an entry.\ 
One may also want to have an index mapping
which indicates at which position the\ 
$i^\mathrm{th}$ entry of the initial list\
ended up getting. In this example, we\  
start out with 
\[\{\lambda_i\}_{i=1}^n=
\{\pgfkeys{/my lists/my initial array/%
content/.print list}\}
\;.\]
The fact that we used the \texttt{/.is array}\ 
key handler means that we know its dimension,\ 
$n=\pgfkeysvalueof{/my lists/my initial array/dim}$.\
After sorting we get 
\[\{\lambda_i'\}_{i=1}^n=\{
\pgfkeys{/my lists/my sorted array/%
content/.print list}\}\]
and an index or reshuffling or permutation list 
\[P=(\pgfkeys{/my lists/%
my index machinery/content/.print list})\;.\]
For instance, the $5^\mathrm{th}$ entry of the\
permutation list being\
$P_5=\pgfkeysvalueof{/my lists/my index machinery/5}$\ 
tells us that the\ 
$\pgfkeysvalueof{/my lists/my index machinery/5}
^\mathrm{th}$\ entry of our original list,\ 
$\lambda_{\pgfkeysvalueof{/my lists/%
my index machinery/5}}
=\pgfkeysvalueof{/my lists/my initial array/%
\pgfkeysvalueof{/my lists/my index machinery/5}}$,\ 
ended up at position $5$ in the sorted list,\ 
i.e.\ $\lambda_5'=\lambda_{\pgfkeysvalueof{/my lists/%
my index machinery/5}}$. Similarly,\ 
$P_3=\pgfkeysvalueof{/my lists/my index machinery/3}$\ 
means that for the\
$\pgfkeysvalueof{/my lists/my index machinery/3}
^\mathrm{th}$ entry of our original list,\  
$\lambda_{\pgfkeysvalueof{/my lists/%
my index machinery/3}}
=\pgfkeysvalueof{/my lists/my initial array/%
\pgfkeysvalueof{/my lists/my index machinery/3}}
=\lambda_3'$. 
\end{minipage}
\end{codeexample}

\subsection{Miscellaneous goodies}

\begin{key}{/tikz/on layer=\meta{layer}}
        Allows one to set a layer in a path. Taken from
		\href{https://tex.stackexchange.com/a/20426}{https://tex.stackexchange.com/a/20426}.
\end{key}

\begin{key}{/tikz/node on layer=\meta{layer}}
        Allows one to set a layer of a node in a path. Taken from
		\href{https://tex.stackexchange.com/a/20426}{https://tex.stackexchange.com/a/20426}.
\end{key}


\begin{key}{/tikz/even odd clip}
        Enables the `even odd rule` for clips.
\end{key}

\begin{key}{/tikz/save named path}
		Saves a path very much like the |save path| key but uses strings, i.e.\
		you can say \texttt{save named path=A}, and prefixes the global macro by
		|tikz@td@named@path@| to minimize the chance that this macro interferes
		with other macros.
\end{key}

\begin{key}{/tikz/use named path}
        Allows one to use a path saved with |save named path|, e.g.\ 
		\texttt{use named path=A}.
\end{key}

\begin{key}{/tikz/same bounding box=\mvar{id} (default A)}
		Allows one to synchronize the bounding box of various pictures. Each picture
		with the same \mvar{id} will have the same bounding box after the second
		compilation.
\end{key}


\subsection{3D-like decorations}

\begin{key}{/tikz/decorations/3d complete coil}
        3d-like coil where the front is thicker than the back.
\end{key}

\begin{key}{/tikz/decorations/3d coil closed}
        Indicates that the coil is closed.
\end{key}

\begin{codeexample}[width=8cm,preamble={\usetikzlibrary{3dtools}}]
\begin{tikzpicture}
\draw[decoration={3d coil color=red,aspect=0.35, segment length=3.1mm, 
amplitude=3mm,3d complete  coil},
decorate] (0,1) -- (0,6);
\draw[decoration={3d coil color=blue,3d coil opacity=0.9,aspect=0.5, 
segment length={2*pi*3cm/50}, amplitude=5mm,3d complete coil,
3d coil closed},
decorate] (5,3.5) circle[radius=3cm];
\end{tikzpicture}
\end{codeexample}


\subsection{Circles of a sphere (compatibility library)}
\label{sec:CircleOfSphere}

|\usetikzlibrary{3dtools.circleofsphere}| provides the commands of the
\href{https://github.com/matthias-wolff/tikz-3dplot-circleofsphere}{%
\texttt{tikz-3dplot-circleofsphere}} package by Matthias Wolff, drawn with
|3d/circle on sphere| and |3d/point on sphere|, so they work with any view and
inside scenes. Load it instead of the package. The styles |tdplotCsFront|,
|tdplotCsBack|, |tdplotCsFill|, |tdplotPtFront| and |tdplotPtBack| keep their
meaning; |tdplotCsDrawAux| is not supported.
|\tdplotCsComputeTransformRotScreen| needs |tikz-3dplot|.

\begin{command}{\tdplotCsDrawCircle\oarg{style}\marg{radius}\marg{alpha}\marg{beta}\marg{elevation}}
	The circle whose plane has the normal with azimuth \meta{alpha} and polar
	angle \meta{beta}, at the elevation angle \meta{elevation} above the center.
	The circle is filled only if \meta{style} mentions |tdplotCsFill|.
\end{command}
\begin{command}{\tdplotCsDrawGreatCircle\oarg{style}\marg{radius}\marg{alpha}\marg{beta}}
\end{command}
\begin{command}{\tdplotCsDrawLatCircle\oarg{style}\marg{radius}\marg{latitude}}
\end{command}
\begin{command}{\tdplotCsDrawLonCircle\oarg{style}\marg{radius}\marg{alpha}}
\end{command}
\begin{command}{\tdplotCsDrawPoint\oarg{style}\marg{radius}\marg{alpha}\marg{beta}}
\end{command}

\begin{codeexample}[width=4.5cm,preamble={\usetikzlibrary{3dtools.circleofsphere}}]
\begin{tikzpicture}[3d/install view={phi=110,theta=70},
  tdplotCsBack/.style={densely dashed,gray}]
 \shade[ball color=white,3d/screen coords,opacity=0.5]
   circle[radius=2];
 \foreach \a in {-60,-30,...,60} {\tdplotCsDrawLatCircle{2}{\a}}
 \foreach \a in {0,30,...,150} {\tdplotCsDrawLonCircle{2}{\a}}
 \tdplotCsDrawPoint{2}{30}{60}
\end{tikzpicture}
\end{codeexample}

\subsection{Remarks on the implementation}

When the library gets loaded, the version of \tikzname\ gets checked. If the
version is earlier than 3.1.1, a warning is issued. In this case, one cannot use
the library unless one updates the pgf installation. Notice that many features
of this library depend on \tikzname's internal macro |\tikz@dcl@coord@|. If this
macro gets removed or changed, many features of the library will no longer work.
The tracking of depths (section~\ref{sec:NamedCoordinates}) wraps
|\pgfpointxyz|, |\pgfpointxy|, |\pgfpointpolarxy|, |\pgfsetxvec| and its
relatives, the |shift| key and several internal macros of \tikzname\ and of its
|calc| library (the current point, named coordinates, node placement, and the
|($...$)| parser). It costs about 1--2\,\% of the run time of pictures that
consist of many plain 2d paths and nodes.
So this whole setup is fragile. Of course, this complication is not unheard of,
but it might be important to know about this before deciding to use this
library. 
% If you have an earlier version
% than 3.1.3, some elaborate pics may not work as they use  |\pgfutil@pushmacro|
% and |\pgfutil@popmacro|, which were added in version 3.1.3. If the version is
% earlier than 3.1.6, the internal macro |\pgfutil@pushmacro| will be fixed.  


\subsection*{Acknowledgments}

I'd like to thank \href{https://github.com/minhthien2016}{minhthien2016} for
continuous support and encouragement. I thank the creators, maintainers,
moderators and users of \href{https://topanswers.xyz/tex}{topanswers.xyz} for
their invaluable support. These acknowledgements extend to all users on this
site, including those with whom I have passionate disagreements. Disagreements
are good, they indicate critical thinking, or in the words of Nobel Laureate
Steven Weinberg, physics thrives on contradictions.


\end{document}


\tdplotsetmaincoords{70}{110} 
\begin{tikzpicture}
 \begin{scope}[local bounding box=tests,tdplot_main_coords]
 % to work with this library, you need to define the cordinate
 % with \path (<x>,<y>,<z>) coordinate (<name>);
  \path (0,0,0) coordinate (O) 
  (1,2,3) coordinate (A) 
  (2,3,-1) coordinate (B) 
  (-1,-2,1) coordinate (C)
  % you can use 3d parse (clumsy)
  [3d parse={0.25*(A)x(B)}] coordinate(D)
  % you can use 3d coordinate to define a new coordinate from existing ones
  [3d coordinate={(E)=0.2*(A)-0.3*(B)+0.6*(C)}] 
  [3d coordinate={(H)=0.2*(A)-0.3*(B)+0.6*(C)}]; 
  \draw (A) -- (B) -- (C) -- (D) -- (E) -- cycle; 
 \end{scope}
 %\RawCoord yields the components
 \edef\tempD{\RawCoord(D)} 
 \edef\tempE{\RawCoord(E)} 
 \edef\tempH{\RawCoord(H)} 
 \node[below right,align=left] at (tests.south west) 
  {$(D)=\tempD$,\\ $(E)=\tempE$,\\ $(H)=\tempH$}; 
\end{tikzpicture} 

\noindent% clumsy parser
$\tdparse{(A)+0.3*(B)>(A)+0.3(B)}=(\pgfmathresult)$

\noindent% parsing inside \pgfmathparse. You need to wrap the argument in "..."
\pgfmathparse{TD("0.2*(A)-0.3*(B)+0.6*(C)")}%
$0.2\,\vec A-0.3\,\vec B+0.6\vec C=(\pgfmathresult)$

%one can parse with the same parser vector products
\noindent\pgfmathparse{TD("0.5*(A)x(B)")}%
$0.5\,\vec A\times\vec B=(\pgfmathresult)$
%(note, however, that something like (A)x(B)x(C) does NOT work)

%as well as scalar products
\noindent\pgfmathparse{TD("(A)+(C)o(B)")}%
$\left(\begin{array}{@{}c@{}}1\\ 0\\ 0\end{array}\right)$
%(note, however, that + and - have higher precedence than o)\end{document}

\end{document}

\endinput
